\documentclass[11pt,a4paper]{amsart}
\usepackage[all]{xy}
\SelectTips{cm}{}
\addtolength{\textwidth}{2cm}
\calclayout

\usepackage{amscd,mathtools,amssymb,amsopn,amsmath,amsthm,graphics,mathrsfs,accents}
\usepackage[colorlinks=true,linkcolor=red,citecolor=blue]{hyperref}
\usepackage[mathscr]{euscript}
\usepackage[usenames]{color}

%\usepackage{showlabels}
%\usepackage{enumerate, pifont, geometry}

%\usepackage{here}

%\usepackage{stmaryrd}
%\usepackage{xcolor}
%\usetikzlibrary{arrows,calc}
%\usepackage{etex}

\newcommand{\com}[1]{{\color{green} #1}}
\newcommand{\old}[1]{{\color{red} #1}}
\newcommand{\new}[1]{{\color{blue} #1}}


\newcommand{\limt}{\underset{\underset{n \in \mathbb{N}}{\longrightarrow}}{\lim}}
\newcommand{\rt}{\rightarrow}
\newcommand{\rat}{\rightarrowtail}
\newcommand{\ri}{\hookrightarrow}
\newcommand{\re}{\twoheadrightarrow}
\newcommand{\lrt}{\longrightarrow}
\newcommand{\lf}{\leftarrow}
\newcommand{\llf}{\longleftarrow}
\newcommand{\ve}{\varepsilon}
\newcommand{\st}{\stackrel}
\def\sw{\operatorname{\mathsf{SW}}}
\newcommand{\Md}{\mathsf{Mod}}
\newcommand{\dl}{\delta}
\def\ds{\operatorname{\mathsf{D_{sg}}}}

\newcommand{\im}{\mathsf{im}}
\def\ker{\operatorname{\mathsf{ker}}}
\def\cok{\operatorname{\mathsf{coker}}}
\newcommand{\fm}{\frak{m}}
\newcommand{\fn}{\frak{n}}
\newcommand{\fp}{\frak{p}}
\newcommand{\fq}{\frak{q}}
\newcommand{\fa}{\frak{a}}
\newcommand{\fb}{\frak{b}}
\newcommand{\T}{\mathcal{T}}

\newcommand{\FF}{\mathbb{F} }
\newcommand{\N}{\mathbb{N} }
\newcommand{\D}{\mathbb{D} }
\newcommand{\Q}{\mathcal{Q} }
\newcommand{\Z}{\mathbb{Z} }
\newcommand{\I}{\mathcal{I} }
\newcommand{\R}{\mathcal{R} }
\newcommand{\C}{\mathscr{C} }

\newcommand{\G}{\mathcal{G} }

\newcommand{\m}{\mathfrak{m}}
\newcommand{\n}{\mathfrak{n}}
\def\qCoh{\operatorname{\mathsf{QCoh}}}
\newcommand{\gl}{\rm{gl} }
\newcommand{\CI}{\mathcal{I} }
\newcommand{\M}{\mathcal{M} }
\newcommand{\Gp}{\rm{Gp} }
\newcommand{\GP}{\rm{GP} }

\newcommand{\GI}{\rm{GI} }

\newcommand{\w}{\mathcal{w} }
\newcommand{\J}{\mathcal{J} }
\newcommand{\CT}{\mathcal{T} }
\def\cp{\operatorname{\mathcal{P}}}

\newcommand{\F}{\mathcal{F} }
\newcommand{\HT}{{\rm{H}}}
\newcommand{\Y}{\mathcal{Y} }
\newcommand{\co}{{\scriptstyle \bf C}\frak{oh}}
\newcommand{\p}{\mathfrak{p} }


\newcommand{\op}{{\rm{op}}}

\newcommand{\End}{{\rm{End}}}
\newcommand{\CM}{{\rm{CM}}}
\newcommand{\depth}{\operatorname{\mathsf{depth}}}
\def\nf{\operatorname{\mathsf{NF}}}
\def\proj{\operatorname{\mathsf{proj}}}
\def\inj{\operatorname{\mathsf{inj}}}
\def\sing{\operatorname{\mathsf{Sing}}}
\def\per{\operatorname{\mathsf{Perf}}}
\def\K{\operatorname{\mathsf{K}}}
\def\Inj{\operatorname{\mathsf{Inj}}}

\newcommand{\ca} {\mathsf{ca}}
\newcommand{\pd} {\mathsf{pd}}
\newcommand{\id} {\mathsf{id}}
\def\db{\operatorname{\mathsf{D^b}}}
\def\thick{\operatorname{\mathsf{thick}}}
\def\ara{\operatorname{\mathsf{ara}}}
\def\fl{\operatorname{\mathsf{fl}}}
\def\Coh{\operatorname{\mathsf{Coh}}}
\def\uHom{\operatorname{\underline{\mathcal{H}{\it om}}}}
\def\Ext{\operatorname{{\mathsf{Ext}}}}

\def\Hom{\operatorname{{\mathcal{H}{\it om}}}}

\def\hom{\operatorname{{\mathsf{Hom}}}}
\def\uhom{\operatorname{\underline{\mathsf{Hom}}}}
\def\eet{\operatorname{{\mathcal{E}{\it xt}}}}
\def\lcm{\operatorname{\mathsf{\underline{CM}}}}
\def\xx{\text{\boldmath $x$}}
\def\yy{\text{\boldmath $y$}}
\def\nf{\operatorname{\mathsf{NF}}}
\def\tot{\operatorname{\mathsf{Tot}}}
\def\spec{\operatorname{\mathsf{spec}}}
\def\sing{\operatorname{\mathsf{Sing}}}
\def\Tor{\operatorname{\mathsf{Tor}}}
\def\hom{\operatorname{\mathsf{Hom}}}
\def\H{\operatorname{\mathsf{H}}}
\def\X{\mathcal{X}}
\def\O{\mathcal{O}}
\def\F{\mathcal{F}}
\def\E{\mathcal{E}}
\def\G{\mathcal{G}}
\def\GG{\mathsf{G}}
\def\Ghd{\mathsf{Ghd}}
\def\hd{\mathsf{hd}}
\def\ext{\operatorname{\mathcal{ext}}}
\def\add{\operatorname{\mathsf{add}}}
\def\md{\operatorname{\mathsf{mod}}}

\newcommand{\uX}{{\mathcal{\underline{X}}}}
\newcommand{\uC}{{\tilde{\mathscr{C}}}}
%\newcommand{\uC}{{\underline{\mathsf{\mathcal{E}xt}}_{\mathscr{{C}}}}}
\def\si{\mathsf{\Sigma}}
\def\syz{\mathsf{\Omega}}
\def\coh{\operatorname{\mathsf{coh}}}
\def\qcoh{\operatorname{\mathsf{Qcoh}}}
\def\GF{\operatorname{\mathsf{GF}}}
\def\supp{\operatorname{\mathsf{Supp}}}
\def\v{\operatorname{\mathsf{V}}}
\def\fin{\operatorname{\mathsf{fin.dim}}}
\def\g{\mathbf g}
\def\a{\mathsf {a}}
\def\x{\mathsf {X}}
\def\k{\mathsf {K}}
\def\pp{\mathsf {P}}
\def\gp{\mathsf {Gp}}
\def\ugp{\underline{\mathsf{Gp}}}

\def\tr{\mathsf {Tr}}
\def\f{\mathsf {F}}
\def\l{\mathsf {L}}
\def\b{\mathsf {b}}
\def\cc{\mathsf {c}}
\def\d{\mathsf {d}}
\def\fl{\mathsf {Flat}}
\def\ch{\mathsf {Ch}}
%\def\f{\mathbf f}
\def\h{\mathbf h}
\def\L{\mathcal L}
\def\p{\mathcal P}
\def\A{\mathcal A}
\def\q{\mathsf Q}
\def\f{\mathbf f}
\def\g{\mathbf g}
\def\u{\mathbf u}
\def\h{\mathbf h}
\def\U{\mathcal U}
\def\co{\operatorname{\mathsf{cone}}}
\def\Proj{\operatorname{\mathsf{Proj}}}
\def\Cot{\operatorname{\mathsf{Cot}}}
\def\ann{\operatorname{\mathsf{ann}}}
\def\ruf{\operatorname{\mathsf{RUF}}}
\def\luf{\operatorname{\mathsf{LUF}}}
\def\c{\operatorname{\underline{\mathscr{C}}}}
\def\Syz{\operatorname{\boldsymbol{\mathsf\Omega}}}
\def\pii{\operatorname{\boldsymbol{\mathbf\pi}}}
\def\phii{\operatorname{\boldsymbol{\mathbf\phi}}}
\def\tth{\operatorname{\boldsymbol{\mathbf\theta}}}
\def\psii{\operatorname{\boldsymbol{\mathbf\psi}}}
\def\vii{\operatorname{\boldsymbol{\mathbf\varphi}}}
\def\vaf{\operatorname{\boldsymbol{\mathbf\varphi}}}

\def\al{\operatorname{\boldsymbol{\alpha}}}
\def\be{\operatorname{\boldsymbol{\beta}}}
\def\ga{\operatorname{\boldsymbol{\gamma}}}
\def\th{\operatorname{\boldsymbol{\theta}}}
\def\ep{\operatorname{\boldsymbol{\epsilon}}}
\def\de{\operatorname{\boldsymbol{\delta}}}
\def\et{\operatorname{\boldsymbol{\eta}}}


\def\ph{\operatorname{\boldsymbol{\phi}}}
\newcommand{\der}{{\rm{der}}}
\newcommand{\rep}{{\rm{rep}}}
\newcommand{\nat}{{\mathsf{Nat}}}

\newcommand{\RR}{\bf{R}}

%\renewcommand{\baselinestretch}{0.27}

%\theoremstyle{plain}
\newtheorem{theorem}{Theorem}[section]
\newtheorem*{mthm}{Main Theorem}
\newtheorem{cor}[theorem]{Corollary}
\newtheorem{properties}[theorem]{Properties}
\newtheorem{lem}[theorem]{Lemma}
\newtheorem{prop}[theorem]{Proposition}
\newtheorem{question}[theorem]{Question}
\newtheorem{conjecture}[theorem]{Conjecture}
\newtheorem{notation}[theorem]{Notation}

\theoremstyle{definition}
\newtheorem{dfn}[theorem]{Definition}
\newtheorem{example}[theorem]{Example}
\newtheorem{examples}[theorem]{Examples}
\newtheorem{conv}[theorem]{Convention}
\newtheorem{construction}[theorem]{Construction}
\newtheorem{obs}[theorem]{Observation}
\newtheorem{rem}[theorem]{Remark}
\newtheorem{s}[theorem]{}

\newtheorem{step}{Step}

\theoremstyle{plain}
\newtheorem{stheorem}{Theorem}[subsection]
\newtheorem{scorollary}[stheorem]{Corollary}
\newtheorem{slemma}[stheorem]{Lemma}
\newtheorem{sproposition}[stheorem]{Proposition}
\newtheorem{squestion}[stheorem]{Question}
\newtheorem{sconjecture}[stheorem]{Conjecture}
\newtheorem{snotation}[stheorem]{Notation}

\theoremstyle{definition}
\newtheorem{sdefinition}[stheorem]{Definition}
\newtheorem{sexample}[stheorem]{Example}
\newtheorem{sobs}[stheorem]{Observation}
\newtheorem*{prob}{Problem}
\newtheorem{sconstruction}[stheorem]{Construction}
\newtheorem{sss}[stheorem]{}
\newtheorem{sremark}[stheorem]{Remark}

\renewcommand{\theequation}{\Alph{equation}}
\numberwithin{equation}{section}
%\usepackage{lineno}
%\linenumbers{lineno}

\begin{document}

\title[On the natural transformations of extension functors]
{On the natural transformations of extension functors}
%\date{today}
\author[Bahlekeh and Salarian]{Abdolnaser Bahlekeh and Shokrollah Salarian}

\address{Department of Mathematics, Gonbad-Kavous University, Postal Code:4971799151, Gonbad-Kavous, Iran}
\email{bahlekeh@gonbad.ac.ir}


\address{Department of Pure Mathematics, Faculty of Mathematics and Statistics,
University of Isfahan, P.O.Box: 81746-73441, Isfahan,
 Iran and \\ School of Mathematics, Institute for Research in Fundamental Science (IPM), P.O.Box: 19395-5746, Tehran, Iran}
 \email{Salarian@ipm.ir}

\subjclass[2020]{18A25, 18G15.}

\keywords{natural transformation; extension functor; Hilton-Rees Theorem.}
\thanks{}


\begin{abstract}
Assume that $\C$ is an exact category.
This paper is concerned with the natural transformations between extension functors on $\C$.  The first main result indicates that if $\C$ has enough projective objects, then for any pair of objects $M, N\in \C$ and any non-negative integer $n$, the group of all natural transformations from $\Ext^{n+1}_{\C}(N, -)$ to $\Ext^{n+1}_{\C}(M, -)$ is isomorphic to the quotient group $\Ext^n_{\C}(M, \syz^nN)/{\p}$, where $\p$ is the subgroup consisting of those extensions of length $n$ arising as a push-out along a morphism $P\rt\syz^nN$, with $P$ projective. This, together with the Auslander-Gruson-Jensen duality yields that if $\C$ is the category of all finitely presented left modules over an associative ring $R$, then the quotient group is isomorphic to the natural transformations from $\Tor_{n+1}^R(-, M)$ to $\Tor_{n+1}^R(-, N)$. The second main result proves  that if $\C$ is an $n$-Frobenuis category, then the statement of the first result remains true, whenever projectives are replaced by $n$-projectives. This result is fruitful from the point of view that, $n$-Frobenius categories may not have projective objects. These results provide a far-reaching generalization of the Hilton-Rees theorem, in the sense that the case $n=0$, recover the Hilton-Rees theorem.
\end{abstract}
\maketitle

\tableofcontents


\section{Introduction}
Assume that $\A$ is an abelian category and $\C$ an extension-closed full additive subcategory of $\A$. So, it will be an exact category in the sense of Quillen, see \cite[Lemma 10.20]{buh}. One of the most important results concerning natural transformations of extension functors on $\C$, is due to  Hilton-Rees \cite{hr}. To be precise, assume that $\C$ has enough projective objects, and $f: M\rt N$ is a morphism in $\C$. Then $\Psi_f:\Ext^1_{\C}(N, -)\lrt\Ext^1_{\C}(M, -)$, in which for any object $T\in\C$ and $\al\in\Ext^1_{\C}(N, T)$, $\Psi_f(T)(\al)=\al f$, is a natural transformation of extension functors, where $\al f$ is the pull-back of $\al$ along $f$. Now, the Hilton-Rees theorem says that the assignment
$$\uhom_{\C}(M, N)\lrt\nat(\Ext^1_{\C}(N, -), \Ext^1_{\C}(M, -)):~~~\bar{f}\mapsto\Psi_f$$is an isomorphism of abelian groups,  see  \cite[Theorem 7]{mart}. Here $\uhom_{\C}(M, N)$ is the group of morphisms from $M$ to $N$ modulo the morphisms factoring through projectives, and  $\nat(\Ext^1_{\C}(N, -), \Ext^1_{\C}(M, -))$ stands for the group of all natural transformations from $\Ext^1_{\C}(N, -)$ to $\Ext^1_{\C}(M, -)$. This, in particular, yields that the latter forms a set. The first aim of this paper is to give a far-reaching generalization of the Hilton-Rees theorem.
Assume that $n$ is a non-negative integer. It is proved in Proposition \ref{1} that, for any pair of objects $N, T\in\C$, and $\al\in\Ext^{n+1}_{\C}(N, T)$, there exists an extension $\alpha_1\in\Ext^1_{\C}(\syz^nN, T)$ and a morphism of extensions $\alpha_1\de_N\lrt\al$, with fixed ends, where $\de_N\in\Ext^n_{\C}(N, \syz^nN)$ is a syzygy sequence. Next, for a given $\ga\in\Ext^n_{\C}(M, \syz^nN)$, we prove that $\Psi_{\ga}:\Ext^{n+1}_{\C}(N, -)\lrt\Ext^{n+1}_{\C}(M, -)$, sending each $\al\in\Ext^{n+1}_{\C}(N, T)$ to $\alpha_1\ga$, is a natural transformation of extension functors, see Corollary \ref{c2}. Now, our first main theorem reads as follows.

\begin{theorem}\label{main1}(Theorem \ref{t2}) Keeping the notation above, the correspondence
$$\Psi:\Ext^n_{\C}(M, \syz^nN)\lrt\nat(\Ext^{n+1}_{\C}(N, -), \Ext^{n+1}_{\C}(M, -)):~~~\ga\mapsto\Psi_{\ga}$$
induces an isomorphism of abelian groups
$$\bar{\Psi}:\Ext^n_{\C}(M, \syz^nN)/{\p}\lrt\nat(\Ext^{n+1}_{\C}(N, -), \Ext^{n+1}_{\C}(M, -)),$$
where $\p$ is the subgroup consisting of those extensions arising as a push-out along a morphism $f: P\rt \syz^nN$ in $\C$, with $P$ projective.
\end{theorem}
It should be noted that  since $\Ext^0_{\C}(M, \syz^0N)=\hom_{\C}(M, N)$, the subgroup $\p$ will be those morphisms $M\rt N$ that factor through some projective objects. This  implies that the case $n=0$ of the above theorem recovers the Hilton-Rees theorem, see Corollary \ref{c33}. Moreover, Theorem \ref{main1}, in conjunction with the Auslander \cite{aus1} and Gruson-Jensen \cite{gj} duality, yields that in the special case when $\C=R$-$\md$, the category of all finitely presented left modules over an associative ring $R$, then the group $\Ext^n_R(M, \syz^nN)/{\p}$ is isomorphic to $\nat(\Tor_{n+1}^R(-, M), \Tor_{n+1}^R(-, N))$, see Corollary \ref{agj}. Particularly, Corollary \ref{vit} reveals that if $\Ext^i_R(M, R)=0$ for all $1\leq i\leq n$, then $\Ext^n_{R}(M, \syz^nN)/{\p}$ is isomorphic to $\uhom_R(M, N)$. These results will be given in Section 3.

The second aim of the paper is to deal with the natural transformations between extension functors over $n$-Frobenius categories. The exact category $\C$ is said to be an $n$-Frobenius category, with $n$ a non-negative integer, if  $\C$ has enough $n$-projective and $n$-injective objects and these two classes of objects coincide, where the notions of $n$-projectivity and $n$-injectivity are defined via the vanishing of $\Ext^{n+1}_{\C}(-, -)$, $(n+1)$-th Yoneda $\Ext$. So, 0-Frobenius categories are exactly Frobenius categories in the usual sense. It is worth noting that $n$-Frobenius categories need not have projective objects, see Example \ref{ex}(3, 6).

Now assume that $\C$ is an $n$-Frobenius category. The second main theorem of the paper asserts that, by replacing projectives with $n$-projectives in Theorem \ref{main1}, the result still remains true. Precisely, we have the result below.

\begin{theorem}\label{main2}(Theorem \ref{nat} and Corollary \ref{ccc}) Let $\C$ be an $n$-Frobenius category and let $M, N$ be objects of $\C$. Then for any  integer $k\geq n$ and a syzygy sequence $\de_N\in\Ext^k_{\C}(N, \syz^kN)$ with respect to $n$-pojectives,  one has   isomorphisms of abelian groups
\[\begin{array}{lllll}
\Ext^k_{\C}(M, \syz^kN)/{\p} & \cong \nat(\Ext^{k+1}_{\C}(N, -), \Ext^{k+1}_{\C}(M, -))\\
& \cong \nat(\Ext^{k+1}_{\C}(-, M), \Ext^{k+1}_{\C}(-, N)).
\end{array}\]
Here $\p$ is the subgroup consisting of those extensions arising as a push-out along a morphism $f:P\rt\syz^kN$, with $P$ $n$-projective.
\end{theorem}
We mention that  the above theorem also provides a vast generalization of  the Hilton-Rees theorem for $n$-Frobenius categories, which need not have projective objects, see Corollary \ref{c44}.  We  illustrate some explicit examples of $n$-Frobenius categories for which the above isomorphisms are satisfied. These results will appear in Section 4.

%\old{ The above theorem,  ensures that for any integer $k>n$,  $\nat(\Ext^{k}_{\C}(N, -), \Ext^{k}_{\C}(M, -))$, as well as, $\nat(\Ext^{k}_{\C}(-, M), \Ext^{k}_{\C}(-, N))$, forms a set. Using this, we provide some conditions in which the natural transformations between (co)homological functors over a ring, forms a set.} 


Throughout the paper, $\C$ is an extension-closed full additive subcategory of an abelian category $\A$.  So, it will be an exact category. We always assume that $n$ is a non-negative integer.
 
%%%%%%%%%%%%%%%%%%%%%%%%%%%%%%%%%%%%%%%%%%%%%%%%%%%%%%%%%%%%%%%%%%%%%%

\section{Preliminaries}

In this section, we recall some notions that will be used in the paper. Particularly, we will remember some facts on extensions briefly and also recall the notion of $n$-Frobenius categories. Moreover, several examples of such categories are provided.

\begin{s}\label{gener}Since $\C$ is an extension-closed subcategory of the abelian category $\A$, the exact structure of $\A$ is inherited by $\C$. Namely, conflations are those short exact sequences in $\A$ with terms in $\C$,  see \cite[Lemma 10.20]{buh}.  We will make use  the following considerations.
\begin{itemize}
\item[(a)] Throughout the paper, instead of conflations, we use the term `extensions'. Also, for simplicity, each extension of the form $0\rt M\rt X\rt N\rt 0$ in $\C$ will be denoted by $M\rt X\rt N$.
\item[(b)] Assume that $n\geq 1$ is an integer. By an extension of length $n$ in $\C$, we mean an exact sequence of a form $\et: 0\rt A\rt X_{n-1}\rt\cdots\rt X_0\rt B\rt 0$ in $\A$, which is obtained by splicing $n$ composable extensions in $\C$. Also,  $A$ and $B$ are called the left and right ends of $\et$, respectively. Such a sequence $\et$ will be denoted by $ A\rt X_{n-1}\rt\cdots\rt X_0\rt B$. If there is no ambiguity, extensions of length $n$ will just be called extensions.
\item[(c)] The set of all equivalence classes of extensions of length $n$ in $\C$, with  $A$ and $B$ being the left and right ends, will be denoted by $\Ext^n_{\C}(B, A)$. We should remind that the equivalence relation is generated by the relation that $\et\sim\et'$, if there is a morphism of extensions $\et\lrt\et'$, with fixed ends. Namely, there is a commutative diagram in $\C$
{\footnotesize\[\xymatrix{\et: A\ar@{=}[d]\ar[r] & X_{n-1}\ar[d]\ar[r] & \cdots\ar[r] & X_0\ar[r]\ar[d] & B\ar@{=}[d]\\ \et': A\ar[r] & X'_{n-1}\ar[r] & \cdots\ar[r] & X'_0\ar[r] & B.}\]}It should be noted that, by the convention, $\Ext^0_{\C}(B, A):=\hom_{\C}(B, A)$. It is known that for any integer $n\geq 1$, $\Ext^n_{\C}(B, A)$, with respect to the Baer sum operation, admits a structure of an abelian group. Indeed, for any pair of elements $\al, \be\in\Ext^n_{\C}(B, A)$, we have $\al+\be=\nabla_A(\al\oplus\be)\Delta_B$. For more details, see \cite[Chapter VII]{mit} and also \cite[Section 9]{hs}. \item[(d)] For given  elements $\et\in\Ext^t_{\C}(B, A)$ and $\et'\in\Ext^{t'}_{\C}(C, B)$, one may splice these together and get  the composition $\et\et'$, which is  an element of $\Ext^{t+t'}_{\C}(C, A)$.
\end{itemize}
\end{s}

\begin{rem}Assume that $f:A\rt A'$ and $g:B\rt B'$ are morphisms in $\C$. Then for any integer $n\geq 0$, the correspondences  $$\Ext^n(B, f):\Ext^n_C(B, A)\lrt\Ext^n_{\C}(B, A') :~~\al\mapsto f\al ~~~and$$
$$\Ext^n(g, A):\Ext^n_{\C}(B', A)\lrt\Ext^n_{\C}(B, A) :~~\be\mapsto \be g$$
are morphisms of abelian groups. Here $f\al$ is the push-out of $\al$ along $f$, and $\be g$ stands for the pull-back of $\be$ along the morphism $g$. One may note that, as $\Ext^0_{\C}(-, -)=\hom_{\C}(-, -)$, these operations can be interpreted as the composition of morphisms in $\C$.
\end{rem}

\begin{theorem}\label{long}Let $\ep: A\st{f}\rt B\st{g}\rt C$ be an extension in $\C$. Then for any object $X\in\C$ and any integer $n>0$, there exists an exact sequence of groups
$$\Ext^{n-1}_{\C}(X, C)\st{\Delta^{n-1}_X}\lrt\Ext^n_{\C}(X, A)\st{\Ext^n(X, f)}\lrt\Ext^n_{\C}(X, B)\st{\Ext^n(X, g)}\lrt\Ext^n_{\C}(X, C)\st{\Delta^n_X}\lrt\Ext^{n+1}_{\C}(X, A),$$where $\Delta^n_{X}(\ga)=\ep\ga$, for any $\ga\in\Ext^n_{\C}(X, C)$.
\end{theorem}
\begin{proof}See \cite[Chapter VII, Theorem 5.1]{mit}.
\end{proof}

\begin{dfn}\label{x}Assume that $\X$ is a class of objects in $\C$ which is closed under finite direct sums. Then for any two objects $A, B\in\C$ and any integer $t\geq 0$, we say that an element $\al\in\Ext^t_{\C}(B, A)$, lies in $\X(B, A)$, if there exists a morphism $f:X\rt A$ in $\C$, with $X\in\X$, and an element $\be\in\Ext^t_{\C}(B, X)$ such that $\al=f\be$. \\
It should be noted that, since $\X$ is closed under finite direct sums, one may easily observe that  $\X(B, A)$ is an additive subgroup of $\Ext^t_{\C}(B, A)$. Moreover, if $t=0$, then $\X(B, A)$ will be the subgroup of $\hom_{\C}(B, A)$ consisting of those morphisms that factor through some objects in $\X$.
\end{dfn}

\begin{dfn}(\cite[Definition 2.3]{bfss}) The exact category $\C$ is called $n$-Frobenius, if
it has enough $n$-projectives and $n$-injectives and these two classes of objects coincide. Recall that
a given object $P\in\C$ (resp. $I\in\C$) is said to be {\it $n$-projective} (resp. {\it $n$-injective}), if $\Ext^i_{\C}(P, X)=0$ (resp. $\Ext^i_{\C}(X, I)=0$) for all integers $i>n$ and all objects $X\in\C$.
\end{dfn}
It follows from the definition that 0-Frobenius categories are exactly Frobenius categories in the usual sense. Also, an $n$-Frobenius category is a $k$-Frobenius cayegory, for any integer $k>n$.  It is known that any abelian category with non-zero $n$-projective objects admits a non-trivial $n$-Frobenius subcategory, see \cite[Theorem 2.15]{bfss}.

\begin{rem}\label{ppp} Keeping the notation of Definition \ref{x}, in this paper, we consider the following two special cases:
\begin{itemize}
\item If $\C$ has enough projective objects, then we take $\X:=\proj\C$, the class of all projective objects of $\C$.
\item If $\C$ is an $n$-Frobenius category, then we set $\X:=n$-$\proj\C$, the class of all $n$-projective objects of $\C$.
\end{itemize}
In both cases, for any pair of objects $ A, B\in\C$, $\X(B, A)$, which will be denoted by $\p(B, A)$, is a subgroup of $\Ext^t_{\C}(B, A)$. If there is no confusion, we denote $\p(-, -)$  by $\p$, and extensions which lie in $\p$ will be called $\p$-extensions.
\end{rem}

Below we give several examples of $n$-Frobenius categories.
\begin{example}\label{ex}
Let $(R, \m)$ be a $d$-dimensional commutative noetherian local ring.\\ (1) The category $\G^{<\infty}$ consisting of all $R$-modules of finite Gorenstein projective dimension is a $d$-Frobenius category, and its $d$-projective objects are indeed the modules of finite projective dimension, see \cite[Example 2.16]{bfss}.\\
{(2) If $R$ is Gorenstein, then the category $\md R$ of all finitely generated $R$-modules, is a $d$-Frobenius category. Particularly, if $d>0$, then injective modules can not be finitely generated, and so, $\md R$ does not have injective objects.\\ (3) Assume that $R$ is  complete  and $\C$ is the class of all artinian $R$-modules. It is known that $D(-):=\hom_{R}(-, E(R/{\m})):\md R\lrt\C$ is a duality. So, according to the previous statement, if $R$ is Gorenstein, then $\C$ will be a $d$-Frobenius category, as well. In particular, if $d>0$, then $\C$ does not have projective objects, because $\md R$ does not have injective objects.}\\ (4) The class $\mathsf{GF}(R)$ of all Gorenstein flat $R$-modules, is a $d$-Frobenius category, and its $d$-projective objects are flat modules, see \cite[Example 2.3(2)]{bfss1}.\\ (5) Assume that $R$ is Gorenstein and $\Q$ a finite acyclic quiver. Then the category  $\G$ of all lattices over $R\Q$ is a 1-Frobenius subcategory of $R\Q$-$\Md$, in which its 1-projective objects are those $R\Q$-modules $M$ which are projective over $R$. We should point out that $\G$ would not be a Frobenius category, if $\Q$ has at least two vertices. Recall that by a lattice over $R\Q$, we mean an $R\Q$-module  which is finitely generated Gorenstein projective over $R$, see \cite[Example 2.6]{bfss1}.\\
(6) Assume that $\coh(\mathbf{P^{1}})$ is the category of  coherent sheaves over the projective line $\mathbf{P^{1}}(R)$, and $\C$
is its full subcategory whose objects are of the form $ \G\equiv G_1 \rt G \lf G_2$, where $G_1$ (resp. $G_2$) is Gorenstein projective over $R[x]$ (resp. $R[x^{-1}]$). Then $\C$ is a 1-Frobenius category which has no projective objects, see \cite[Appendix A]{bfss1} and \cite[Corollary 2.3]{ee}.
\end{example}

\begin{s}\label{ss}Assume that $\C$ is an $n$-Frobenius category.\\
(1) Since $\C$ has enough $n$-projective objects, for any object $M\in\C$ and an integer $t\geq 1$, one may obtain a syzygy sequence  $\syz^tM\rt P_{t-1}\rt\cdots\rt P_0\rt M$, where each $P_i$ is  $n$-projective.  Also, $\syz^tM$ is called a $t$-th syzygy of $M$. Syzygy sequences are also called unit extensions.\\ (2)  For a given object $X\in\C$, the class of all  unit  extensions (of length $n$)  ending at  $X$ (resp.  beginning with  $X$) is denoted by $\U_n(X)$ (resp.  {$\U^n(X)$}).\\ (2) Assume that $t\geq 1$ is an integer and $\delta, \delta' \in \U^t(N)$ (resp. $\delta_1, \delta_2 \in\U_t(N)$). Then there exist deflations (resp. inflations) $a, a'$ with $n$-projective kernels (resp. cokernels) such that $\delta a=\delta'a'$ (resp. $a\delta_1=a'\delta_2$), see \cite[Proposition 5.5]{bfss}. In this situation, $[\delta a, \delta'a']$  (resp. $[a\delta_1, a'\delta_2]$) is called a co-angled pair (resp. an angled pair). In some cases, we denote it by $\delta\st{a}\llf\delta''\st{a'}\lrt\delta'$ (resp. $\delta_1\st{a}\lrt\delta''\st{a'}\llf\delta_2$), where $\delta''=\delta a$  (resp. $\delta''=a\delta_1$).\\ (3) Suppose that an element $\ga\in\Ext^t_{\C}(M, \syz^tN)$ is given, where $t\geq 1$ is an integer. So there exist unit extensions $\delta\in\U^t(\syz^t N)$, $\delta'\in\U_t(M)$ and morphisms $f,g$ in $\C$ such that $\ga=\delta f$ and $\ga=g\delta'$, see \cite[Proposition 5.2]{bfss}. These are called a right and a left unit factorization of $\ga$, respectively, and abbreviated by $\ruf$ and $\luf$.
\end{s}


\begin{s}{\sc $n$-$\Ext$-phantom (invertible) morphisms.} Assume that $\C$ is an $n$-Frobenius category and $f:M\rt N$ is a morphism in $\C$.\\ (1) For a given object $X\in\C$, there exist induced maps $\Ext^n_{\C}(X, M)/{\p}\rt\Ext^n_{\C}(X, N)/{\p}$ and $\Ext^n_{\C}(N, X)/{\p}\rt\Ext^n_{\C}(M, X)/{\p}$ that assign  each object $\ga+\p$ to ${f\ga}+\p$ and ${\ga f}+\p$, respectively.\\ (2)
$f$ is called an $n$-$\Ext$-invertible morphism, provided that   for any object $X\in\C$, the induced morphism $\Ext^n_{\C}(X, M)/{\p}\rt\Ext^n_{\C}(X, N)/{\p}$ is an isomorphism. It can be seen that,  this is equivalent to saying that $\Ext^n_{\C}(N, X)/{\p}\rt\Ext^n_{\C}(M, X)/{\p}$) is also an isomorphism,  see \cite[Corollaries 3.4 and 4.8]{bfss}. These, in particular, imply that $f$ is an $n$-$\Ext$-invertible morphism if and only if the natural transformation $\Ext^{n+1}(-, f):\Ext^{n+1}_{\C}(-,M)\lrt\Ext^{n+1}_{\C}(-,N)$ (or equivalently, $\Ext^{n+1}(f, -)$) is an equivalence of functors. 
It should be pointed out that $n$-$\Ext$-invertible morphisms are called quasi-invertible morphisms in \cite{bfss}. Such morphisms are also called quasi-isomorphisms in \cite{hr}.\\
 (3) $f$ is said to be an $n$-$\Ext$-phantom morphism, if   for any object $X\in\C$, the induced morphism $\Ext^n_{\C}(X, M)/{\p}\rt\Ext^n_{\C}(X, N)/{\p}$  (or equivalently, $\Ext^n_{\C}(N, X)/{\p}\lrt\Ext^n_{\C}(M, X)/{\p}$) is zero, see \cite[Theorem 6.7 and Definition 6.8]{bfss}. By \cite[Corollary 6.9]{bfss},  $f$ is an $n$-$\Ext$-phantom morphism, if and only if  the natural transformation $\Ext^{n+1}(-, f):\Ext^{n+1}_{\C}(-,M)\lrt\Ext^{n+1}_{\C}(-,N)$ (or equivalently, $\Ext^{n+1}(f, -)$) vanishes.  
\end{s}

\begin{example}(1) Over an $n$-Frobenius category $\C$, any morphism
with $n$-projective kernel and cokernel is an  $n$-$\Ext$-invertible morphism, see \cite[Corollary 3.6]{bfss}. Also, a morphism that factors through an $n$-projective object is also an $n$-$\Ext$-phantom morphism.\\ (2)  Over Frobenius categories (or equivalently, 0-Frobenius categories), 0-$\Ext$-phantom morphisms are those morphisms that factor through projective objects, see \cite[Proposition 3.8(1)]{bfss1}.\\ (3) If $\C$ is a Frobenius category, then a given morphism $f: M\rt N$ in $\C$ is a 0-$\Ext$-invertible if  and only if there exist $P,Q\in\proj\C$ such that $M\oplus Q\st{{{\tiny {\left[\begin{array}{ll} f & g_1 \\ l & g_2 \end{array} \right]}}}}\lrt N\oplus P$ is an isomorphism, see \cite[Proposition 3.8(2)]{bfss1}.
\end{example}

\begin{rem}\label{invert}Assume that $\C$ is an $n$-Frobenius category and  an element $\ga\in\Ext^n_{\C}(M, N)$ is given. It is known that for any $n$-$\Ext$-invertible morphism $a:M'\rt M$, $\ga$ is a $\p$-extension if and only if $\ga a$ is a $\p$-extension. Similarly, $b\ga$ is a $\p$-extension if and only if $\ga$ is so, whenever $b:N\rt N'$ is an $n$-$\Ext$-invertible morphism, see \cite[Proposition 4.9]{bfss}.
\end{rem}

\section{A generalization of the Hilton-Rees theorem}
The purpose of this section is to provide a far-reaching generalization of the Hilton-Rees theorem. Namely, we give the proof of  Theorem \ref{main1}, stated in the introduction.

Throughout this section, we always assume that $\C$ has enough projective objects and $n$ is a non-negative integer. Let us begin with an auxiliary result.

\begin{prop}\label{1}Let $\al$ be an element of $\Ext^{n+1}_{\C}(N, T)$. Then for a given syzygy sequence $\de_N\in\Ext^n_{\C}(N, \syz^nN)$, there exists an element $\alpha_1\in\Ext^1_{\C}(\syz^nN, T)$ such that there is a morphism $\alpha_1\de_N\lrt\al$ with fixed ends.
\end{prop}
\begin{proof}Decompose $\al: T\rt X_n\rt\cdots\rt X_0\rt N$ into two extensions $\alpha':T\rt X_n\rt L$ and $\alpha'': L\rt X_{n-1}\rt\cdots\rt X_0\rt N$. Now assuming $\de_N:\syz^nN\rt P_{n-1}\rt\cdots\rt P_0\rt N$, where each $P_i$ is projective, one may obtain the following commutative diagram with exact rows: {\footnotesize\[\xymatrix{\de_N:\syz^{n}N\ar[d]_{f}\ar[r] & P_{n-1}\ar[d]\ar[r] & \cdots\ar[r] & P_0\ar[r]\ar[d] & N\ar@{=}[d]\\
\alpha'':L\ar[r] & X_{n-1}\ar[r] & \cdots\ar[r] & X_0\ar[r] & N.}\]}Now take a pull-back diagram
{\footnotesize \[\xymatrix{\alpha_1:T\ar[r]\ar@{=}[d] & G\ar[r]\ar[d]_{f_1} & \syz^{n}N\ar[d]_{f} \\ \alpha':T\ar[r] & X_{n}\ar[r]& L.}\] }
Splicing the above diagrams together gives us the following commutative diagram:
{\footnotesize\[\xymatrix{\alpha_1\de_N:T\ar[r]\ar@{=}[d]& G\ar[d]\ar[r] & P_{n-1}\ar[d]\ar[r] & \cdots\ar[r] & P_0\ar[r]\ar[d] & N\ar@{=}[d]\\ \al: T\ar[r]& X_{n}\ar[r] & X_{n-1}\ar[r] & \cdots\ar[r] & X_0\ar[r] & N,}\]}as desired.
\end{proof}

\begin{rem}\label{r1}Assume that $\al\in\Ext^{n+1}_{\C}(N, T)$ is given. Take a syzygy sequence $\de_N\in\Ext^n_{\C}(N, \syz^nN)$.
According to Proposition \ref{1}, there exists a morphism $\alpha_1\de_N\lrt\al$ with fixed ends, where $\alpha_1\in\Ext^1_{\C}(\syz^nN, T)$. Now for a given element $\ga\in\Ext^n_{\C}(M, \syz^nN)$, $\alpha_1\ga$ will be an element of $\Ext^{n+1}_{\C}(M, T)$.
\end{rem}

\begin{theorem}\label{t1}With the notation of Remark \ref{r1}, the assignment
$$\Psi_{\ga}(T): \Ext^{n+1}_{\C}(N, T)\lrt\Ext^{n+1}_{\C}(M, T):~~\al\mapsto \alpha_1\ga $$
is a morphism of abelian groups.
\end{theorem}
\begin{proof}
The first issue to address is well-definedness. In this direction, assume that there exists an element  $\alpha'_1\in\Ext^1_{\C}(\syz^nN, T)$ and a morphism with fixed ends $\alpha'_1\de_N\lrt\al$. We shall prove that $\alpha_1\ga=\alpha'_1\ga$. By our construction, we have the following commutative diagram:
{\footnotesize\[\xymatrix{\de_N:\syz^{n}N\ar[d]_{f'}\ar[r] & P_{n-1}\ar[d]\ar[r] & \cdots\ar[r] & P_0\ar[r]\ar[d] & N\ar@{=}[d]\\ \alpha'':L\ar[r] & X_{n-1}\ar[r] & \cdots\ar[r] & X_0\ar[r] & N.}\]}Now Comparison Theorem (cf.  \cite[Theorem 6.16]{rot}), ensures that $f-f'$ factors through a projective object, and so, taking the following pull-back diagram
{\footnotesize \[\xymatrix{T\ar[r]\ar@{=}[d] & G''\ar[r]\ar[d] & {\syz}^{n}N\ar[d]_{f-f'} \\ \alpha':T\ar[r] & X_n\ar[r]& L,}\]}we infer that the upper row, $\alpha_1-\alpha'_1,$ splits. This, in turn, yields that $\alpha_1\ga=\alpha'_1\ga$, as desired. Next, we show that $\Psi_{\ga}(T)$ is a morphism of abelian groups. So, take elements  $\al, \be\in\Ext^{n+1}_{\C}(N, T)$. It should be proved that $\Psi_{\ga}(T)(\al)+\Psi_{\ga}(T)(\be)=\Psi_{\ga}(T)(\al+\be)$. According to Proposition \ref{1}, there are $\alpha_1, \beta_1\in\Ext^1_{\C}(\syz^nN, T)$ and morphisms of extensions $\alpha_1\de_N\lrt\al$ and $\beta_1\de_N\lrt\be$ with fixed ends. So $\alpha_1\de_N+\beta_1\de_N\lrt\al+\be$ is a morphism of extensions with fixed ends. On the other hand, by \cite[Chapter VII, Lemma 3.2]{mit},  the equality $\alpha_1\de_N+\beta_1\de_N=(\alpha_1+\beta_1)\de_N$ holds true. Thus, by our definition, one obtains the following equalities:
$$\Psi_{\ga}(T)(\al+\be)=(\alpha_1+\beta_1)\ga=\alpha_1\ga+\beta_1\ga=\Psi_{\ga}(T)(\al)+\Psi_{\ga}(T)(\be).$$
So the proof is completed.
\end{proof}

\begin{cor}\label{c2}Let  $\de_N\in\Ext^n_{\C}(N, \syz^nN)$  be a syzygy sequence. Then for any element $\ga\in\Ext^n_{\C}(M, \syz^nN)$, $\Psi_{\ga}(-):\Ext^{n+1}_{\C}(N, -)\lrt\Ext^{n+1}_{\C}(M, -)$ is a natural transformation of extension functors.
\end{cor}
\begin{proof}In view of Theorem \ref{t1}, for any object $T\in\C$, $\Psi_{\ga}(T):\Ext^{n+1}_{\C}(N, T)\lrt\Ext^{n+1}_{\C}(M, T)$ is a morphism of abelian groups. Moreover, for any morphism $g:T\rt T'$ in $\C$, one evidently has the following commutative square
\[\xymatrix{\Ext^{n+1}_{\C}(N, T)\ar[r]^{\Psi_{\ga}(T)}\ar[d]_{\Ext^{n+1}(N,g)}&\Ext^{n+1}_{\C}(M, T)\ar[d]_{\Ext^{n+1}(M, g)}\\ \Ext^{n+1}_{\C}(N, T')\ar[r]^{\Psi_{\ga}(T')}&\Ext^{n+1}_{\C}(M, T').}\]Consequently, $\Psi_{\ga}(-)$ is a natural transformation from  $\Ext^{n+1}_{\C}(N, -)$ to $\Ext^{n+1}_{\C}(M, -)$.
It should be noted that, since $\Ext^0_{\C}(M, \syz^0N)=\hom_{\C}(M, N)$, in the case $n=0$, $\ga:M\rt N$ will be a  morphism in $\C$, and also, for a given element  $\al\in\Ext^1_{\C}(N, T)$, $\al=\alpha_1$. In particular, $\Psi_{\ga}=\Ext^1(\ga, -):\Ext^{1}_{\C}(N, -)\lrt\Ext^{1}_{\C}(M, -)$ is a well-known natural transformation.
So the proof is finished.
\end{proof}

\begin{cor}\label{c3}Let $\de_N\in\Ext^n_{\C}(N, \syz^nN)$ be a syzygy sequence. Then for any object $M\in\C$, the corrspondence $$\Psi: \Ext^{n}_{\C}(M, \syz^nN)\lrt \nat(\Ext^{n+1}_{\C}(N, -), \Ext^{n+1}_{\C}(M, -)):~~\ga\mapsto \Psi_{\ga} $$ is a morphism of abelian groups.
\end{cor}
\begin{proof}
Assume that an extension $\ga\in\Ext^n_{\C}(M, \syz^nN)$ is given. By Corollary \ref{c2}, $\Psi_{\ga}$ lies in $\nat(\Ext^{n+1}_{\C}(N, -), \Ext^{n+1}_{\C}(M, -))$. Suppose that $\ga=\ga'$, as elements of $\Ext^n_{\C}(M, \syz^nN)$. We shall prove that $\Psi_{\ga}=\Psi_{\ga'}$. As mentioned in \ref{gener}(c), we may assume that, there exists a morphism of extensions $\ga\lrt\ga'$ with fixed ends. Take an arbitrary elemen $\al\in\Ext^{n+1}_{\C}(N, T)$. In view of Proposiion \ref{1}, there is a morphism $\alpha_1\de_N\lrt\al$. Thus, one may obtain a morphism of extensions $\alpha_1\ga\lrt\alpha_1\ga'$ with fixed ends, and so, $\alpha_1\ga=\alpha_1\ga'$, as objects of $\Ext^{n+1}(M, T)$, see \cite[Chapter VII, Proposition 3.1]{mit}. This means that $\Psi_{\ga}(T)(\al)=\Psi_{\ga'}(T)(\al),$ which  ensures the well-defindness of $\Psi$. Next take two elements $\ga, \ga''\in\Ext^n_{\C}(M, \syz^nN)$. So one has the equalities $$\Psi_{\ga+\ga''}(T)(\al)=\alpha_1(\ga+\ga'')=\alpha_1\ga+\alpha_1\ga''=\Psi_{\ga}(T)(\al)+\Psi_{\ga''}(T)(\al),$$ implying that $\Psi$ is a morphism of abelian groups. Hence the proof is completed.
\end{proof}

Before giving the proof of the main result of this section, we include a couple of axuillary results. 

Recall that for any pair of objects $X, Y\in\C$, an element $\ga\in\Ext^n_{\C}(X, Y)$ lies in $\p$, if there is a morphism $f:P\rt Y$ in $\C$, with $P$ projective, and an element $\et\in\Ext^n_{\C}(X, P)$ such that $\ga=f\et$. As mentioned before, $\p$ is a subgroup of $\Ext^n_{\C}(X, Y)$, see Definition \ref{x} and Remark \ref{ppp}.

 \begin{lem}\label{p} Let $\ep: X\rt Y\st{h}\rt Z$ be an extension in $\C$ and let $\Delta_M^n: \Ext^n_{\C}(M, Z)\rt\Ext^{n+1}_{\C}(M, X)$ be the connecting homomrphism. If $\ga\in\Ext^n_{\C}(M, Z)$ lies in $\p$, then $\Delta^n_M(\ga)=0$. Moreover, the converse is true, whenever $Y$ is a projective object of $\C$.
\end{lem}
\begin{proof}
Since $\ga\in\p$, by the definition, there exists an extension $\be\in\Ext^n_{\C}(M, P)$ and a morphism $f:P\rt Z$  in $\C$ with $P\in\proj\C$, such that $\ga=f\be$. As $P$ is projective, $\ep f$ is split, and so, $(\ep f)\be$ is a zero element of $\Ext^{n+1}_{\C}(M, X).$  By \cite[Chapter VII, Proposition 3.1]{mit},  $(\ep f)\be=\ep\ga$. Now since $\Delta^n_M(\ga)=\ep\ga$, the claim follows. For the converse, assume that $\Delta^n_M(\ga)=0$. We would like to show that $\ga\in\p$.  By Theorem \ref{long}, there exists an exact sequence of abelian groups: $$\Ext^n_{\C}(M, Y)\st{\Ext^n(M, h)}\lrt\Ext^n_{\C}(M, Z)\st{\Delta^n_M}\lrt\Ext^{n+1}_{\C}(M, X).$$Since $\Delta_M^n(\ga)=0$, one finds an element  $\al\in\Ext^n_{\C}(M, Y)$ such that $h\al=\ga$, gives the desired result, because $Y\in\proj\C$. So the proof is finished.
\end{proof}

The result below reveals that any natural transformation $\Psi$ between extension functors, is determined by its action  on syzygy sequences.

\begin{prop}\label{epi1}Let  $\Phi:\Ext^{n+1}_{\C}(N, -)\rt\Ext^{n+1}_{\C}(M, -)$ be a natural transformation, and  $\de_N\in\Ext^{n+1}_{\C}(N, \syz^{n+1}N)$  an arbitrary syzygy sequence. Then for any $T\in\C$ and  $\al\in\Ext^{n+1}_{\C}(N, T)$, there exists a morphism $f:\syz^{n+1}N\rt T$ in $\C$ such that $\Phi_T(\al)=\Ext^{n+1}(M, f)\Phi_{\syz^{n+1}N}(\de_N)$.
\end{prop}
\begin{proof}
Consider a syzygy sequence $\de_N:\syz^{n+1}N\rt P_n\rt\cdots\rt P_0\rt N$.
Assume that an element $\al\in\Ext^{n+1}_{\C}(N, T)$ is given. Since each $P_i$ is projective, one may obtain the
following commutative diagram
{\footnotesize\[\xymatrix{\de_N:\syz^{n+1}N\ar[d]_{f}\ar[r] & P_{n}\ar[d]\ar[r] & \cdots\ar[r] & P_0\ar[r]\ar[d] & N\ar@{=}[d]\\ \al:T\ar[r] & X_n\ar[r]  & \cdots\ar[r] & X_0\ar[r] & N.}\]}As explained in \cite[Remark 2.7(1)]{bfss}, we may assume that $\al=f\de_N$. Consider the following commutative square:
\[\xymatrix{\Ext^{n+1}_{\C}(N, \syz^{n+1}N)\ar[r]^{\Phi_{\syz^{n+1}N}}\ar[d]_{\Ext^{n+1}(N, f)} & \Ext^{n+1}_{\C}(M, \syz^{n+1}N)\ar[d]_{\Ext^{n+1}(M, f)} \\ \Ext^{n+1}_{\C}(N, T)\ar[r]^{\Phi_T} & \Ext^{n+1}_{\C}(M, T).}\]So, we have the following equalities:  $$\Phi_T(\al)=\Phi_T\Ext^{n+1}(N, f)(\de_N)=\Ext^{n+1}(M, f)\Phi_{\syz^{n+1}N}(\de_N),$$as desired. So the proof is completed.
\end{proof}

Now we are ready to give the proof of the main result of this section.
\begin{theorem}\label{t2} Keeping the notation of Corollary \ref{c3}, the correspondence
$$\bar{\Psi}: \Ext^n_{\C}(M, \syz^nN)/{\p}\lrt\nat(\Ext^{n+1}_{\C}(N, -), \Ext^{n+1}_{\C}(M, -)):~~\bar{\ga}\mapsto \Psi_{\ga}$$is an isomorphism of abelian groups.
\end{theorem}
\begin{proof}
By Corollary \ref{c3}, $\Psi:\Ext^n_{\C}(M, \syz^nN)\lrt\nat(\Ext^{n+1}_{\C}(N, -), \Ext^{n+1}_{\C}(M, -))$, sending each $\ga$ to $\Psi_{\ga}$, is a morphism of abelian groups.  We claim that $\Psi$ is an epimorphism. To see this, take a natural transformation $\Phi:\Ext^{n+1}_{\C}(N,-)\lrt\Ext^{n+1}_{\C}(M, -)$. By making use of Proposition \ref{epi1}, we only need to show that for a given syzygy sequence $\de_N\in\Ext^{n+1}_{\C}(N, \syz^{n+1}N)$, there exists an element $\ga\in\Ext^n_{\C}(M, \syz^nN)$ such that  $\Psi_{\ga}(\de_N)=\Phi_{\syz^{n+1}N}(\de_N)$. Decompose $\de_N=\delta\de'_N$ with $\delta:\syz^{n+1}N\st{a}\rt P_n\rt\syz^nN$ and $\de'_N\in\Ext^n_{\C}(N, \syz^nN)$. So one obtains the following commutative diagram of abelian groups with exact rows:
\[\xymatrix{\Ext^n_{\C}(N, \syz^nN)\ar[r]^{\Delta^n_N} & \Ext^{n+1}_{\C}(N, \syz^{n+1}N)\ar[r]^{\epsilon}\ar[d]_{\Phi_{\syz^{n+1}N}}& \Ext^{n+1}_{\C}(N, P_n)\ar[d]_{\Phi_{P_n}} \\ \Ext^n_{\C}(M, \syz^nN)\ar[r]^{\Delta^n_M}& \Ext^{n+1}_{\C}(M, \syz^{n+1}N)\ar[r]^{\eta}& \Ext^{n+1}_{\C}(M, P_n) .}\]Since $\epsilon\Delta_N^n(\de'_N)=\epsilon(\delta\de'_N)=\epsilon(\de_N)=0$, we have $\eta\Phi_{\syz^{n+1}N}(\de_N)=0.$ Thus there is an element $\ga\in\Ext^n_{\C}(M, \syz^nN)$ such that $\Delta^n_M(\ga)=\Phi_{\syz^{n+1}N}(\de_N)$, namely, $\Phi_{\syz^{n+1}N}(\de_N)=\delta\ga$. By our definition, $\delta\ga=\Psi_{\ga}(\syz^{n+1}N)(\de_N)$, and so,  $\Phi_{\syz^{n+1}N}(\de_N)=\Psi_{\ga}(\syz^{n+1}N)(\de_N)$, as desired. Hence, it remains to show that $\ker(\Psi)=\p$.  First, one should note that by Theorem \ref{t1}, for any element $\ga\in\Ext^n_{\C}(M, \syz^nN)$ and $\al\in\Ext^{n+1}_{\C}(N, T)$, we have $\Psi_{\ga}(T)(\al)=\Delta_M^n(\ga)$, where $\Delta_M^n:\Ext^n_{\C}(M, \syz^nN)\lrt\Ext^{n+1}_{\C}(M, T)$ is the connecting homomorphism associated to $\alpha_1:T\rt G\rt\syz^nN$. This fact, in conjunction with  Lemma \ref{p}, ensures that each object in $\p$ lies in $\ker(\Psi)$. Next,  assume that $\ga\in\Ext^n_{\C}(M, \syz^nN)$ such that $\Psi_{\ga}=0$. So, for a given syzygy sequence $\de_N\in\Ext^{n+1}_{\C}(N, \syz^{n+1}N)$, $0=\Psi_{\ga}(\syz^{n+1}N)(\de_N)=\Delta_M^n(\ga)$, where $\Delta_M^n:\Ext^n_{\C}(M, \syz^nN)\lrt\Ext^{n+1}_{\C}(M, \syz^{n+1}N)$ is the connecting homomorphism associated to the extension $\syz^{n+1}N\rt P_n\rt\syz^nN$. Now the reverse implication of Lemma \ref{p} yields that $\ga\in\p$. So the proof is completed.
\end{proof}

The result below indicates that, in the case $n=0$, Theorem \ref{t2} recovers the Hilton-Rees theorem, see \cite[Theorem 7]{mart}.
\begin{cor}\label{c33}For any pair of objects $M, N\in\C$, the correspondence
$$\uhom_{\C}(M, N)\lrt\nat(\Ext^1_{\C}(N, -), \Ext^1_{\C}(M, -)):~~\bar{f}\mapsto\Ext^1(f, -)$$
is an isomorphism of abelian groups.
\end{cor}
\begin{proof}Since $\Ext^0_{\C}(M, \syz^0N)=\hom_{\C}(M, N)$, and the push-out is the composition of morphisms,  the subgroup $\p$ of  $\Ext^0_{\C}(M, \syz^0N)$, exactly contains those morphisms $M\rt N$ that factor through a projective object. That is, $\Ext^0_{\C}(M, \syz^0N)/{\p}=\uhom_{\C}(M, N)$. Now Theorem \ref{t2} finishes the proof.
\end{proof}

In the following, we will observe that if $\C=R$-$\md$, the category of all finitely presented left $R$-modules, then the quotient group in Theorem \ref{t2} can be specified as natural transformations between $\Tor$-functors. In this direction,  the Auslander-Gruson-Jensen duality plays a crucial role.  So let us recall it from \cite{aus1, gj}.

\begin{s}{\sc Auslander--Gruson-Jensen duality}\\   Assume that $R$ is an associative ring with identity, and $R$-$\md$ (resp. $\md$-$R$) is the category of finitely presented left (resp. right) $R$-modules. Assume that ${\mathsf{fp}}(R$-$\md, \mathsf{Ab})$ is the category consisting of all finitely presented additive functors $F:R$-$\md\lrt\mathsf{Ab}$, where $\mathsf{Ab}$ is the category of abelian groups. It is known that this is an abelian category with kernels and cokernels computed object-wise, see for example \cite[Section 2]{aus}. The Auslander-Gruson-Jensen duality, which was first written down explicitly by Auslander in \cite{aus1} and by Gruson
and Jensen in \cite{gj}, says that, there is an equivalence
\begin{center}
$d: {\mathsf{fp}}(R$-$\md, \mathsf{Ab})\lrt ({\mathsf{fp}}(\md$-$R, \mathsf{Ab}))^{op},$
\end{center}
sending any finitely presented functor $F:R$-$\md\rt\mathsf{Ab}$ to the functor $dF:\md$-$R\rt\mathsf{Ab}$, defined by $(dF)M=\nat(F, M\otimes_R-)$ for any finitely presented right $R$-module $M$. In particular, if $F=\hom_R(X, -)$, for some $X\in R$-$\md$, then by using Yoneda Lemma, we infer that $d(\hom_R(X, -))=-\otimes_RX$.
\end{s}
The Auslander--Gruson-Jensen duality, combined with Theorem \ref{t2}, yields the following interesting result.
\begin{cor}\label{agj}For any  finitely presented left $R$-modules $M, N$, there is an isomorphism of groups
$$\Ext^n_{\C}(M, \syz^nN)/{\p}\cong\nat(\Tor_{n+1}^R(-, M), \Tor_{n+1}^R(-, N)).$$
\end{cor}
\begin{proof}
Assume that $X$ and $Y$ are finitely presented left and right $R$-module, respectively.  Since $Y\otimes_R-$ is a right exact functor, by making use of \cite[Theorem 1.2]{hr} or \cite[Corollary 5.4]{aus}, one obtains the following isomorphism: $$\nat(\Ext^{n+1}_R(X, -), Y\otimes_R-)\cong L_{n+1}(Y\otimes_R-)(X)=\Tor_{n+1}^R(Y, X),$$which is natural in $X$ and $Y$. This, in particular, means that $d(\Ext^{n+1}_R(X, -))=\Tor_{n+1}^R(-, X)$, which is natural in $X$. Now the Auslander--Gruson-Jensen duality, together with Theorem \ref{t2} yields the desired result.
\end{proof}

Let us close this section with a couple of results. The authors thank Vitor  Gulisz for pointing these out to them.
\begin{cor}\label{vit}Let $M, N$ be finitely presented left $R$-modules. If $\Ext^i_R(M, R)=0$ for all $1\leq i\leq n$, then $\Ext^n_R(M, \syz^nN)/{\p}\cong\uhom_{R}(M, N)$.
\end{cor}
\begin{proof}Consider the following isomorphisms:
\[\begin{array}{lllll}
\Ext^n_R(M, \syz^nN)/{\p} & \cong \nat(\Tor^{R}_{n+1}(-, M), \Tor^{R}_{n+1}(-, N))\\
& \cong \nat(\Tor^{R}_1(-, \syz^nM), \Tor^R_1(-, \syz^nN))\\ &\cong\uhom_R(\syz^nM, \syz^nN)\\
&\cong\Tor_1^R(\tr\syz^nM, \syz^nN)\\
&\cong\Tor_{n+1}^R(\tr\syz^nM, N)\\
&\cong\uhom_R(M, N),
\end{array}\]
where $\tr\syz^nM$ denotes the transpose of $\syz^nM$. The first isomorphism follows from Corollary \ref{agj}, while the second and fifth ones hold true, thanks to the dimension shifting argument. The validity of the third isomorphism comes from \cite[Lemma 9]{mart}, and the fourth one is well-known, see for example \cite[Lemma 3.9]{yo} or \cite[Proposition 65]{gu}. Finally, the last isomorphism follows from \cite[Proposition 1.1.3]{iy}. So the proof is finished.
\end{proof}

Let $\Lambda$ be an artin algebra, that is, it contains a central artinian subring $R$ and $\Lambda$ is a finitely generated $R$-module. Then $D(-):=\hom_R(-, E(R/{J_R})): \Lambda$-$\md\lrt\md$-$\Lambda$ is a duality, where $J_R$ is the Jacobson radical of $R$. For any integer $n\geq 1$, $\tau_n(-):=D\tr\syz^n(-)$, is called the $n$-Auslander-Reiten translation functor. This notion has been defined by Iyama \cite{iy}, when he developed his
higher Auslander-Reiten theory.

Now \cite[Theorem 1.5]{iy} combined with Corollary \ref{vit}, immediately yields the following result.
\begin{cor}Let $\Lambda$ be an artin algebra and let $M, N$ be finitely generated $\Lambda$-modules. If $\Ext^i_{\Lambda}(M, \Lambda)=0$ for all $1\leq i\leq n$, then $\Ext^n_{\Lambda}(M, \syz^nN)/{\p}\cong D\Ext^{n+1}_{\Lambda}(N, \tau_{n+1}M).$
\end{cor}

\section{Natural transformations of extension functors over $n$-Frobenius categories }
This section contains the proof of Theorem \ref{main2}, which concerns the natural transformations between extension functors over an $n$-Frobenius category. By the aid of this, we provide some situations in which natural transformations of (co)homology functors over a ring form a set. Let us  begin  by recalling the following useful remark.

\begin{rem}\label{lr}Assume that $\C$ is an $n$-Frobenius category and $X, Y$ are two objects of $\C$. Then a given extension $\ga\in\Ext^n_{\C}(X, Y)$ lies in $\p$, if there is a morphism $f:P\rt Y$ in $\C$, where $P$ is $n$-projective, and an extension $\be\in\Ext^n_{\C}(X, P)$ such that $\ga=f\be$, see Remark \ref{ppp}. It is worthwhile to note that since $\C$ is $n$-Frobenius, \cite[Proposition 4.3]{bfss} indicates that $\ga\in\p$ if and only if there exists a morphism $g:X\rt Q$ in $\C$ with $Q$ $n$-projective, and $\et\in\Ext^n_{\C}(Q, Y)$ such that $\ga=\et g$.
\end{rem}

\begin{lem}\label{pp}Let $\C$ be an $n$-Frobenius category and let
$\ga\in\Ext^n_{\C}(X, Y)$ be a $\p$-extension. Then for any extension $\alpha\in\Ext^1_{\C}(Y, Z)$, the composition $\alpha\ga=0$ in $\Ext^{n+1}_{\C}(X, Z)$.
\end{lem}
\begin{proof}Since $\ga$ is a $\p$-extension, as mentioned in Remark \ref{lr}, there exists a morphism $h:X\rt P$ in $\C$, where $P$ is $n$-projective, and an element $\be\in\Ext^n_{\C}(P, Y)$ such that $\ga=\be h$. As $P$ is $n$-projective, $\Ext^{n+1}_{\C}(P, -)=0$. So, for a given extension $\alpha\in\Ext^1_{\C}(Y, Z)$, we have that $\alpha\be=0$, implying that $(\alpha\be)h=\alpha(\be h)=\alpha\ga=0$, as desired.
\end{proof}

\begin{rem}Assume that $\C$ has enough projective objects. Then the lifting property yields that for any object $N\in\C$ and an integer $n\geq 1$, the $n$-th syzygy $\syz^nN$ of $N$ is unique, up to projective summands. However, this is not the case for $n$-Frobenius categories, where $\syz^nN$ is obtained using $n$-projective objects. Our key strategy to fix this flaw is the use of (co)angled pairs. Indeed, as observed in \ref{ss}(2), for any pair of syzygy sequence $\de_{N}, \de'_N\in\U_n(N)$, there is an angled pair $\de_N\st{a}\rt\de''_N\st{b}\lf\de'_{N}$. In particular, there are morphisms $\syz^nN\st{a}\rt{\syz''}^nN\st{b}\lf{\syz'}^nN$, which are inflations with $n$-projective cokernels, and so, they are $n$-$\Ext$-invertible morphisms.
\end{rem}

In the remainder of the paper, we always assume that $\C$ is an $n$-Frobenus category, with $n$ a non-negative integer.

\begin{prop}\label{ll}For a given element  $\al\in\Ext^{n+1}_{\C}(N, T)$, there exists a unit extension $\de_N\in\Ext^n_{\C}(N, \syz^nN)$ and an extension $\alpha_1\in\Ext^1_{\C}(\syz^nN, T)$ such that there is a morphism $\alpha_1\de_N\rt\al$ with fixed ends. Moreover, if $\be\in\Ext^{n+1}_{\C}(N, T)$ is another extension, then there exists a unit extension $\de'_N\in\Ext^n_{\C}(N, {\syz'}^nN)$ and morphisms $\alpha_1\de'_N\rt\al$ and $\beta_1\de'_N\rt\be$ with fixed ends.
\end{prop}
\begin{proof}Decompose $\al=\alpha_n\al'$ with $\alpha_n: T\rt X_n\rt L$ and $\al': L\rt X_{n-1}\rt\cdots\rt X_0\rt N$. In view of \cite[Lemma 5.2(2)]{bfss},   one has the following commutative diagram with exact rows:
{\footnotesize\[\xymatrix{\de_N:{\syz}^{n}N\ar[d]_{f}\ar[r] & P_{n-1}\ar[d]\ar[r] & \cdots\ar[r] & P_0\ar[r]\ar[d] & N\ar@{=}[d]\\ \al':L\ar[r] & X_{n-1}\ar[r] & \cdots\ar[r] & X_0\ar[r] & N,}\]}where each $P_i$ is $n$-projective. Setting $\alpha_1:=\alpha_nf$, one obtains the morphism of extensions $\alpha_1\de_N\lrt\al$ with fixed ends. For the second assertion, decompose $\be=\beta_n\be'$ with $\beta_n:T\rt Y_n\rt L'$ and $\be': L'\rt Y_{n-1}\rt\cdots\rt Y_0\rt N$. Set $\et:=(\al'\oplus\be')\Delta$, where $\Delta:N\st{{{\tiny {\left[\begin{array}{ll} 1 \\ 1 \end{array} \right]}}}}\lrt N\oplus N$ is a morphism in $\C$. Another use of \cite[Lemma 5.2(2)]{bfss}, yields the following commutative diagram:
{\footnotesize\[\xymatrix{\de'_N:{\syz'}^{n}N\ar[d]_{{{{\tiny {\left[\begin{array}{ll} g \\ g' \end{array} \right]}}}}}\ar[r] & P'_{n-1}\ar[d]\ar[r] & \cdots\ar[r] & P'_0\ar[r]\ar[d] & N\ar@{=}[d]\\ \et:L\oplus L'\ar[r] & T_{n-1}\ar[r] & \cdots\ar[r] & T_0\ar[r] & N.}\]}Now putting $\alpha_1:=\alpha_ng$ and $\beta_1:=\beta_ng'$, one may obtain the morphism of extensions $\alpha_1\de'_N\lrt\al$ and $\beta_1\de'_N\lrt\be$ with fixed ends. So the proof is finished.
\end{proof}

\begin{s}Assume that $a:M\rt M'$ is an $n$-$\Ext$-invertible morphism in $\C$ and $\ga\in\Ext^n_{\C}(M, N)$ is given. So  one may find an element $\ga'\in\Ext^n_{\C}(M', N)$ such that $\ga-\ga'a$ is a $\p$-extension. To simplify the notation, we put $\ga':=_{_{\p}}\ga a^{-1}$. Similarly, for any element $\be\in\Ext^n_{\C}(N, M')$, there is  $\be'\in\Ext^n_{\C}(N, M)$ such that $\be -a\be'$ is a $\p$-extension. In this case, we write $\be':=_{_{\p}}a^{-1}\be$.
\end{s}

\begin{rem}\label{3s1}Assume that a unit extension $\de_N\in\Ext^n_{\C}(N, \syz^nN)$ and  $\al\in\Ext^{n+1}_{\C}(N, T)$ are given. According to Proposition \ref{ll}, there exists a unit extension $\de'_N\in\Ext^{n}_{\C}(N, {\syz'}^nN)$ and a morphism of extensions $\alpha_1\de'_N\lrt\al$ with fixed ends, in which $\alpha_1\in\Ext^1_{\C}({\syz'}^nN, T)$.
Now for a given element $\ga\in\Ext^n_{\C}(M, \syz^nN)$, we infer that $b^{-1}a\ga\in\Ext^n_{\C}(M, {\syz'}^nN)$, where $\de_N\st{a}\rt\de''_N\st{b}\lf\de'_N$ is an angled pair. In particular, setting $\ga':=b^{-1}a\ga$, we have that $\alpha_1\ga'\in\Ext^{n+1}_{\C}(M, T)$.
\end{rem}
\begin{prop}\label{3p1}With the notation of  Remark \ref{3s1}, $$\Psi_{\ga}(T):\Ext^{n+1}_{\C}(N, T)\lrt\Ext^{n+1}_{\C}(M, T):~~ \al\mapsto\alpha_1\ga'$$ is a morphism of abelian groups.
\end{prop}
\begin{proof}First we prove that $\Psi_{\ga}(T)$ is well-defined. To do this, assume that there exists a unit extension $\de_{1N}\in\Ext^n_{\C}(N, \syz^n_1N)$ and a morphism $f':\syz^n_1N\rt L$ in which $(\alpha_nf')\de_{1N}\lrt\al$ is a morphism of extensions with fixed ends. Now considering an angled pair $\de_N\st{a_1}\rt\de_1\st{b_1}\lf\de_{1N}$, we prove that $(\alpha_nf)b^{-1}a\ga=_{\p}(\alpha_nf')b_1^{-1}a_1\ga$. One should note that, since there exists a morphism of extensions $(\alpha_nf)b^{-1}a\ga\lrt\alpha_n(f(b^{-1}a\ga))$ with fixed ends, by \cite[Chapter VII, Proposition 3.1]{mit}, $(\alpha_nf)b^{-1}a\ga=\alpha_n(f(b^{-1}a\ga))$. Analoguesly,  $(\alpha_nf')b_1^{-1}a_1\ga=\alpha_n(f'(b_1^{-1}a_1\ga))$. So it suffices to show that  $\alpha_n(f(b^{-1}a\ga))=\alpha_n(f'(b_1^{-1}a_1\ga))$. To this end, similar to the proof of  \cite[Theorem 7.4]{bfss}, one obtains the following diagram of angled pairs:
{\footnotesize \[\xymatrix{&\delta_N\ar[dl]_{a}\ar[dr]^{a_1}& \\ \delta''_{N}~\ar[r]^{a_3}& \delta_3& \delta_1\ar[l]_{b_3}\\ \delta'_{N}\ar[u]^{b}\ar[r]^{a_2}~ & \delta_2\ar[u]^{c_3} & \delta_{1N}\ar[l]_{b_2}\ar[u]^{b_1}.}\]}By making use of \cite[Proposition 5.8]{bfss}, there is a morphism $h$ in $\C$ such that $\ga=h\de_2$, $f=ha_2$ and $f'=hb_2$. In particular, one may have the following diagram in $\C$: {\footnotesize\[\xymatrix{\de'_N~ \ar[r]^{a_2}\ar[rd]_{f}& \de_2\ar[d]^{h}& {\de_1}_{N}\ar[l]_{b_2}\ar[ld]^{f'}\\ & \ga .& }\]}
Thus, we have the following equalities modulo $\p$: $$f'(b_1^{-1}(a_1\ga))=_{\p}hb_2(b_1^{-1}(a_1\ga))=_{\p}h(b_2(b_1^{-1}(a_1\ga)))=_{\p}fa_2^{-1}(b_2(b_1^{-1}(a_1\ga)))$$ $$=_{\p}f(a_2^{-1}(b_2(b_1^{-1}(a_1\ga))))=_{\p} fa_2^{-1}(c_3^{-1}b_3)(a_1\ga)))=_{\p}f(a_2^{-1}c_3^{-1})(b_3(a_1\ga))$$$$=_{\p}f(b^{-1}a_3^{-1})(b_3(a_1\ga))=_{\p}f(b^{-1}(a_3^{-1}b_3a_1)\ga)=_{\p}f(b^{-1}(a\ga)).$$Now Lemma \ref{pp} together with Remark \ref{invert}, ensures that $\Psi_{\ga}(T)$ is well-defined.
%Next assume that $\ga$ is a $\p$-conflation. We should prove that $\Psi_{\ga}(T)(\al)=\alpha_1(b^{-1}a\ga)=0$. Since $\ga$ is a $\p$-conflation, by \cite[Proposition 4.9]{bfss}, $\ga'=b^{-1}a\ga$ is a $\p$-conflations, as well. Consequently, Lemma \ref{pp} ensures that $\alpha_1\ga'=0$, and so, $\Psi_{\ga}(T)$ is well-defined.
Finally, we show that $\Psi_{\ga}(T)$ is a morphism of abelian groups. To this end, consider extensions $\al, \be\in\Ext^{n+1}_{\C}(N, T)$. In view of Proposition \ref{ll}, there are morphisms of extensions $\alpha_1\de'_N\lrt\al$ and $\beta_1\de'_N\lrt\be$ with fixed ends, for some unit extension $\de'_N\in\Ext^n_{\C}(N, {\syz'}^nN)$. This, in conjunction with \cite[Chapter VII, Lemma 3.2]{mit}, gives us the morphism of extensions $(\alpha_1+\beta_1)\de'_N\lrt\al+\be$ with fixed ends. Now putting $\ga':=b^{-1}a\ga$, where $\de_N\st{a}\rt\de''_N\st{b}\lf\de'_N$ is an angled pair, one may obtain the following equalities:
$$\Psi_{\ga}(T)(\al+\be)=(\alpha_1+\beta_1)\ga'=\alpha_1\ga'+\beta_1\ga'=\Psi_{\ga}(T)(\al)+\Psi_{\ga}(T)(\be).$$Thus the proof is completed.
\end{proof}

The result below is proven similar to that of Corollary \ref{c3}. However, for the reader's convenience, its proof is included.
\begin{prop}\label{p33}Let $\de_N\in\Ext^n_{\C}(N, \syz^nN)$ be a unit extension. Then for any object $M\in\C$, the correspondence
$$\Psi:\Ext^n_{\C}(M, \syz^nN)\lrt\nat(\Ext^{n+1}_{\C}(N, -), \Ext^{n+1}_{\C}(M, -)) ~~: ~~\ga\mapsto \Psi_{\ga}$$
is a morphism of abelian groups. 
\end{prop}
\begin{proof}First we show that $\Psi$ is well-defined. In this direction, take an object $T\in\C$. By  Proposition \ref{3p1}, $\Psi_{\ga}(T):\Ext^{n+1}_{\C}(N, T)\lrt\Ext^{n+1}_{\C}(M, T)$ is a morphism of abelian groups. In particular, one may follow the argument given in the proof of Corollary \ref{c2} and deduce that $\Psi_{\ga}:\Ext^{n+1}_{\C}(N, -)\lrt\Ext^{n+1}_{\C}(M, -)$ is a natural transformation. Next assume that $\ga$ and $\ga'$ are two elements of $\Ext^n_{\C}(M, \syz^nN)$ such that there exists a morphism of extensions $\ga\lrt\ga'$ with fixed ends. For a given object $T\in\C$, we shall prove that $\Psi_{\ga}(T)=\Psi_{\ga'}(T)$. In this direction, take an arbitrary element $\al\in\Ext^{n+1}_{\C}(N, T)$. In view of Proposition \ref{ll}, there exists a unit extension $\de'_N\in\Ext^n_{\C}(N, {\syz'}^nN)$ and $\alpha_1\in\Ext^1_{\C}({\syz'}^nN, T)$, in which there is a morphism of extensions $\alpha_1\de'_N\lrt\al$ with fixed ends. Consider an angled pair $\de_N\st{a}\rt\de''_N\st{b}\lf\de'_N$. Since there is a morphism of extensions $\ga\lrt\ga'$ with fixed ends, using the universal property of push-out, one may find a morphism of extensions $b^{-1}a\ga\lrt b^{-1}a\ga'$ with fixed ends. Thus, $\alpha_1(b^{-1}a\ga)\lrt\alpha_1( b^{-1}a\ga')$ will also be a morphism of extensions with fixed ends. Consequently, as elements of $\Ext^{n+1}_{\C}(M, T)$, these extensions are the same, meaning that  $\Psi_{\ga}(T)(\al)=\Psi_{\ga'}(T)(\al)$, as desired. Thus $\Psi$ is well-defined. Next, we show that $\Psi$ is a morphism of groups. To this end, considering a pair of extensions $\ga, \ga''\in\Ext^n_{\C}(M, \syz^nN)$, one gets the following equalities: $$\Psi_{\ga+\ga''}(T)(\al)=\alpha_1(b^{-1}a(\ga+\ga''))=\alpha_1(b^{-1}a\ga)+b^{-1}a\ga'')=\alpha_1(b^{-1}a\ga)+\alpha_1(b^{-1}a\ga'').$$ But the right-hand side equals to $\Psi_{\ga}(T)(\al)+\Psi_{\ga''}(T)(\al)$, and so,  the proof is finished.
\end{proof}


% The following remark indicates that, by replacing projectives with $n$-projectives in Lemma \ref{p}, the result still remains true for $n$-Frobenius categories.
%\begin{rem}\label{4.1} Let $\ep: X\rt Y\rt Z$ be an extension in $\C$. For a given object $M\in\C$, consider the connecting homomorphism $\Delta_M^n:\Ext^n(M, Z)\lrt\Ext^{n+1}(M, X).$ Assume that a conflation $\ga\in\Ext^n(M, Z)$ lies in $\p$. Since $\C$ is an $n$-Frobenius category, in view of \cite[Theorem 6.7]{bfss}, there exists a conflation $\be\in\Ext^n(P, Z)$ and a morphism $f:M\rt P$, where $P$ is $n$-projective, such that $\ga=\be f$. So one has the equalities $$\Delta_M^n(\ga)=\ep\ga=\ep(\be f)=(\ep\be)f.$$ Now since $P$ is $n$-projective, we infer that $\Ext^{n+1}(P, X)=0$, and in particular, $\ep\be=0$. Consequently, $\Delta_M^n(\ga)=(\ep\be)f=0$. Next assume that in extension $\ep$, $Y$ is $n$-projective, and $\Delta_M^n(\ga)=0$. Then the same argument given in the proof of reverse implication in Lemma \ref{p}, yields that $\ga\in\p$.  \end{rem}

\begin{prop}\label{epi}Let  $\Phi:\Ext^{n+1}_{\C}(N, -)\rt\Ext^{n+1}_{\C}(M, -)$ be a natural transformation and  $\de_N\in\Ext^{n+1}_{\C}(N, \syz^{n+1}N)$ be a unit extension. Then for any object $T\in\C$ and  $\al\in\Ext^{n+1}_{\C}(N, T)$, there exists a morphism $f:{\syz'}^{n+1}N\rt T$ in $\C$ and $n$-$\Ext$-invertible morphisms  $\syz^{n+1}N\st{a}\rt{\syz''}^{n+1}N\st{b}\lf{\syz'}^{n+1}N$ such that $\Phi_T(\al)=f^*{b^*}^{-1}a^*\Phi_{\syz^{n+1}N}(\de_N)$, where $(-)^*=\Ext^{n+1}(M, -)$.
\end{prop}
\begin{proof}
Assume that $T\in\C$ and $\al\in\Ext^{n+1}_{\C}(N, T)$ is given. Take an $\luf$ $\al=f\de'_N$ of $\al$, where $f:{\syz'}^{n+1}N\rt T$ is a morphism in $\C$. Consider an angled pair $\de_N\st{a}\rt\de''_N\st{b}\lf\de'_N$. Now we obtain the following commutative square:
\[\xymatrix{\Ext^{n+1}_{\C}(N, \syz^{n+1}N)\ar[r]^{\Phi_{\syz^{n+1}N}}\ar[d]_{f^*{b^*}^{-1}a^*} & \Ext^{n+1}_{\C}(M, \syz^{n+1}N)\ar[d]_{f^*{b^*}^{-1}a^*} \\ \Ext^{n+1}_{\C}(N, T)\ar[r]^{\Phi_T} & \Ext^{n+1}_{\C}(M, T).}\]
One should note that $(-)^*$ in the left (resp. right) column stands for $\Ext^{n+1}(N, -)$ (resp. $\Ext^{n+1}(M, -)$). So, going in the counterclockwise side, one gets the equalities $\Phi_T f^*{b^*}^{-1}a^*(\de_N)=\Phi_Tf^*(\de'_N)=\Phi_T(\al)$. Now since the square is commutative, we have that $\Phi_T(\al)=f^*{b^*}^{-1}a^*\Phi_{\syz^{n+1}N}(\de_N)$, where $(-)^*=\Ext^{n+1}(M, -)$. So the proof is finished.
\end{proof}

%\begin{rem}\label{4.3}Assume that $M, N\in\C$ and $\de_N\in\Ext^n(N, \syz^nN)$ is a unit conflation. For a given element $\ga\in\Ext^n(M, \syz^nN)$, we assign a natural transformation $\Phi_{\ga}:\Ext^{n+1}(N, -)\lrt\Ext^{n+1}(M, -)$, as follows. Assume that $T\in\C$ and  $\al: T\rt X_n\rt\cdots\rt X_0\rt N$ is an arbitrary object of $\Ext^{n+1}(N, T)$. Consider an angled pair $\de_N\st{a}\rt\de''_N\st{b}\lf\de'_N$ and set $\ga':=b^{-1}a\ga$. Now define $\Psi_{\ga}(\al):=\Delta_M^n(f\ga')$, where $\Delta_M^n:\Ext^n(M, {\syz'}^nN)\lrt\Ext^{n+1}(M, T)$, the connecting homomorphism associated to the conflation $\alpha_1: T\rt G\rt{\syz'}^nN$.\end{rem}

Now we are ready to prove the main result of this section.

\begin{theorem}\label{nat} Let $M, N$ be arbitrary objects of $\C$ and let $\de_N\in\Ext^n_{\C}(N, \syz^nN)$ be a unit extension. Then the correspondence
$$\Psi: \Ext^n_{\C}(M, \syz^nN)\lrt{\nat}(\Ext^{n+1}_{\C}(N, -), \Ext^{n+1}_{\C}(M, -))~:~~\ga\mapsto \Psi_{\ga}$$
induces an isomorphism
$$\Ext^n_{\C}(M, \syz^nN)/{\p}\st{\cong}\lrt{\nat}(\Ext^{n+1}_{\C}(N, -), \Ext^{n+1}_{\C}(M, -)).$$
\end{theorem}
\begin{proof}
As observed in Proposition \ref{p33}, for any object $M\in\C$,  $$\Psi:\Ext^n_{\C}(M, \syz^nN)\lrt\nat(\Ext^{n+1}_{\C}(N,-), \Ext^{n+1}_{\C}(M, -))~:~~\ga\mapsto\Psi_{\ga}$$ is a morphism of abelian groups. We would like to show that $\Psi$ is an epimorphism. In this direction, take a natural transformation $\Phi:\Ext^{n+1}_{\C}(N,-)\lrt\Ext^{n+1}_{\C}(M, -)$. Consider a unit extension $\be=\delta\de_N\in\Ext^{n+1}_{\C}(N, \syz^{n+1}N)$ with $\delta:\syz^{n+1}N\rt P_n\rt\syz^nN$. So one obtains the following commutative diagram with exact rows:\[\xymatrix{\Ext^n_{\C}(N, \syz^nN)\ar[r]^{\Delta^n_N} & \Ext^{n+1}_{\C}(N, \syz^{n+1}N)\ar[r]^{\epsilon}\ar[d]_{\Phi_{\syz^{n+1}N}}& \Ext^{n+1}_{\C}(N, P_n)\ar[d]_{\Phi_{P_n}} \\ \Ext^n_{\C}(M, \syz^nN)\ar[r]^{\Delta^n_M}& \Ext^{n+1}_{\C}(M, \syz^{n+1}N)\ar[r]^{\eta}& \Ext^{n+1}_{\C}(M, P_n),}\]where $\Delta^n:\Ext^n_{\C}(-, \syz^nN)\lrt\Ext^{n+1}_{\C}(-, \syz^{n+1}N)$ is the connecting homomorphism associated to $\delta:\syz^{n+1}N\rt P_n\rt\syz^nN$. {One should note that since $P_n$ is $n$-projective,  it is an $n$-injective object, implying that $\Ext^{n+1}_{\C}(-, P_n)=0$. This, in particular, yields that $\Delta^n_M$ is an epimorphism.  Hence, }  one may find an element $\ga\in\Ext^n_{\C}(M, \syz^nN)$ such that $\Delta_M^n(\ga)=\Phi_{\syz^{n+1}N}(\be)$. On the other hand, since $\delta\de_N\st{=}\lrt\be$ is a morphism with fixed ends, by our definition, $\Psi_{\ga}(\syz^{n+1}N)(\be)=\delta\ga=\Delta_M^n(\ga)$. Consequently, $\Psi_{\ga}(\syz^{n+1}N)(\be)=\Phi_{\syz^{n+1}N}(\be).$ Now Proposition \ref{epi} ensures that $\Psi$ is an epimorphism. Hence, in order to complete the proof, we need to show that $\ker\Psi=\p$. To this end, assume that $\ga\in\Ext^n_{\C}(M, \syz^nN)$ is a $\p$-extension. Then, by making use of  Lemma \ref{pp} and Remark \ref{invert}, one may infer that $\Psi_{\ga}=0$. Conversely, assume that $\ga\in\ker\Psi$. So, taking $\al:=\delta\de_N$, with $\delta:\syz^{n+1}N\rt P_n\rt\syz^nN$, we have $\Psi_{\ga}(\syz^{n+1}N)(\al)=\delta\ga=0$. Thus, $\Delta^n_M(\ga)=\delta\ga=0$, where $\Delta_M^n:\Ext^n_{\C}(M, \syz^nN)\lrt\Ext^{n+1}_{\C}(M, \syz^{n+1}N)$ is the connceting homomorphism associted to $\delta$. Now, since $P_n$ is $n$-projective,  the same argument given in the proof of the reverse implication in Lemma \ref{p}, yields that $\ga\in\p$. Hence, the proof is completed.
\end{proof}

\begin{rem}\label{rr}Assume that $M, T$ are objects in $\C$ and $\al\in\Ext^{n+1}_{\C}(T, M)$ is given. Decompose $\al=\al'\alpha'_1$ with $\al':M\rt X_n\rt\cdots\rt X_1\rt L$ and $\alpha'_1:L\rt X_0\rt T$. In view of \cite[Lemma 5.2(1)]{bfss},   one has the following commutative diagram with exact rows:
{\footnotesize\[\xymatrix{\al':~~M\ar@{=}[d]\ar[r] & X_n\ar[d]\ar[r] & \cdots\ar[r] & X_1\ar[r]\ar[d] & L\ar[d]_{f}\\ \de_{\syz^{-n}M}:M\ar[r] & P_{n}\ar[r] & \cdots\ar[r] & P_1\ar[r] & \syz^{-n}M,}\]}where each $P_i$ is $n$-projective. Now, putting $\alpha_1:=f\alpha'_1$, one obtains the morphism of extensions $\al\lrt\de_{\syz^{-n}M}\alpha_1$ with fixed ends.
\end{rem}

\begin{s}\label{cov}Assume that $M, N$ are objects of $\C$ and $\de_{\syz^{-n}M}\in\Ext^n_{\C}(\syz^{-n}M, M)$ is a unit extension. Then for a given element $\ga\in\Ext^n_{\C}(\syz^{-n}M, N)$, we assign a natural transformation $\Phi_{\ga}:\Ext^{n+1}_{\C}(-, M)\lrt\Ext^{n+1}_{\C}(-, N)$,  as follows:\\ For a given object $T\in\C$, we take an element $\al\in\Ext^{n+1}_{\C}(T, M)$. As mentioned in Remark \ref{rr}, there is a morphism of extensions $\al\lrt\de_{{\syz'}^{-n}M}\alpha_1$ with fixed ends, where $\de_{{\syz'}^{-n}M}\in\Ext^n_{\C}({\syz'}^{-n}M, M)$ is a unit extension and $\alpha_1\in\Ext^1_{\C}(T, {\syz'}^{-n}M)$. So considering a co-angled pair $\de_{\syz^{-n}M}\st{a}\lf\de_{{\syz''}^{-n}M}\st{b}\rt\de_{{\syz'}^{-n}M}$, one has $(\ga a)b^{-1}\in\Ext^n_{\C}({\syz'}^{-n}M, N)$. Now we define $\Phi_{\ga}(T)(\al):=((\ga a)b^{-1})\alpha_1$, which lies in $\Ext^{n+1}_{\C}(T, N)$. A dual manner given in the proof of Proposition \ref{p33} and using Remark \ref{lr}, reveals that $\Phi_{\ga}:\Ext^{n+1}_{\C}(-, M)\lrt\Ext^{n+1}_{\C}(-, N)$ is a natural transformation.  In particular, as $\C$ has enough $n$-injectives and $n$-injectives coincide with $n$-projectives,   dualizing the argument, yields the result below.
\end{s}

\begin{theorem}\label{th3}With the notation above, the correspondence
$$\Phi:\Ext^n_{\C}(\syz^{-n}M, N)\lrt\nat(\Ext^{n+1}_{\C}(-, M), \Ext^{n+1}_{\C}(-, N))~:\ga\mapsto\Phi_{\ga}$$ induces an isomorphism of abelian groups
$$\bar{\Phi}:\Ext^n_{\C}(\syz^{-n}M, N)/{\p}\st{\cong}\lrt\nat(\Ext^{n+1}_{\C}(-, M), \Ext^{n+1}_{\C}(-, N))~:\bar{\ga}\mapsto\Phi_{\ga}.$$ 
\end{theorem}

\begin{cor}\label{cc}Let $M, N$ be objects of $\C$ and $\de_N\in\Ext^n_{\C}(N, \syz^nN)$ a unit extension. Then there is an isomorphism $$\Ext^n_{\C}(M, \syz^nN)/{\p}\cong\nat(\Ext^{n+1}_{\C}(-, M), \Ext^{n+1}_{\C}(-, N)).$$
\end{cor}
\begin{proof}
In view of Theorem \ref{th3}, it suffices to show that for a given unit extension $\de_{\syz^{-n}M}\in\Ext^n_{\C}(\syz^{-n}M, M)$, an isomorphism  $\Ext^n_{\C}(\syz^{-n}M, N)/{\p}\cong\Ext^n_{\C}(M, \syz^nN)/{\p}$ holds true. To do this, assume that $\de_N:=\syz^nN\rt P_{n-1}\rt\cdots\rt P_0\rt N$. Since $P_0$ is $n$-projective, and so, it is an $n$-injective object, $\Ext^{n+1}_{\C}(-, P_0)=0$. Thus, we give an exact sequence $$\Ext^n_{\C}(\syz^{-n}M, P_0)\lrt\Ext^n_{\C}(\syz^{-n}M, N)\lrt\Ext^{n+1}_{\C}(\syz^{-n}M, \syz N)\lrt 0.$$This yields an isomorphism $\Ext^n_{\C}(\syz^{-n}M, N)/{\p}\cong\Ext^{n+1}_{\C}(\syz^{-n}M, \syz N)$. Similarly, considering the unit extension $\syz^{-n+1}M\rt Q_n\rt\syz^{-n}M$, and applying Remark \ref{lr}, one gets an isomorphism $\Ext^n_{\C}(\syz^{-n+1}M, \syz N)/{\p}\cong \Ext^{n+1}_{\C}(\syz^{-n}M, \syz N)$. Combining the latter two isomorphisms would give us the isomorphism of groups $\Ext^n_{\C}(\syz^{-n}M, N)/{\p}\cong\Ext^n_{\C}(\syz^{-n+1}M, \syz N)/{\p}$. Now one may continue this argument repeatedly, and deduce that $\Ext^n_{\C}(\syz^{-n}M, N)/{\p}\cong\Ext^n_{\C}(M, \syz^n N)/{\p}$, as desired.
\end{proof}

Since $\C$ is an $n$-Frobenius category, as noted in \cite[Remark 2.4]{bfss}, it is a $k$-Frobenius category with $n$-$\proj\C=k$-$\proj\C$, for all integers $k>n$. So, as a direct consequence of Theorem \ref{nat} and Corollary \ref{cc}, we include the following result.
\begin{cor}\label{ccc}Let $M, N$ be objects of $\C$ and $k\geq n$ an integer, and let $\de_N\in\Ext^k_{\C}(N, \syz^kN)$ be a unit extension. Then one has the following isomorphisms of abelian groups
\[\begin{array}{lllll}
\Ext^k_{\C}(M, \syz^kN)/{\p} & \cong \nat(\Ext^{k+1}_{\C}(N, -), \Ext^{k+1}_{\C}(M, -))\\
& \cong \nat(\Ext^{k+1}_{\C}(-, M), \Ext^{k+1}_{\C}(-, N)).
\end{array}\]
\end{cor}
The above corollary immediately yields the following interesting result.
\begin{cor}For any pair of objects $M, N\in\C$ and an integer $k\geq n$,  the groups $\nat(\Ext^{k+1}_{\C}(M, -), \Ext^{k+1}_{\C}(N, -))$ and $\nat(\Ext^{k+1}_{\C}(-, M), \Ext^{k+1}_{\C}(-, N))$ form a set.
\end{cor}


Since over Frobenius categories (or equivalently, 0-Frobenius categories), $\p$-extensions are those morphisms that factor through a projective object, the case $n=0$ of Theorem \ref{nat} and Corollary \ref{cc} recovers Theorems 7 and 8 of \cite{mart}, which is the Hilton-Rees theorem for Frobenius categories. Indeed, we have the result below.

\begin{cor}\label{c44}Let $\C$ be a Frobenius category. Then  for any two objects $M, N\in\C$, one has the following isomorphisms of abelian groups
\[\begin{array}{lllll}
\underline{\hom}_{\C}(M, N) & \cong \nat(\Ext^{1}_{\C}(N, -), \Ext^{1}_{\C}(M, -))\\
& \cong \nat(\Ext^{1}_{\C}(-, M), \Ext^{1}_{\C}(-, N)).
\end{array}\]
\end{cor}
We close this paper by presenting several interesting examples of $n$-Frobenius categories that satisfy the  isomorphisms appeared in Corollary \ref{ccc}.

\begin{example}Let $(R, \m)$ be a commutative Gorenstein local ring with $\dim R=d$.\\
(1) Since by \cite[Theorem 10.2.4]{ej}, $d$-th syzygy of any $R$-module is Gorenstein projective, one may see that the category $R$-$\Md$ of all $R$-modules is a $d$-Frobenius category, where $d$-projectives are modules of finite projective dimension. So, for any pair of $R$-modules $M, N$ and $k\geq d$, the following isomorphisms of groups
\[\begin{array}{lllll}
\Ext^k_R(M, \syz^kN)/{\p} & \cong \nat(\Ext^{k+1}_R(N, -), \Ext^{k+1}_R(M, -))\\
& \cong \nat(\Ext^{k+1}_R(-, M), \Ext^{k+1}_R(-, N)),
\end{array}\]
hold. We should emphasize that $\syz^kN$ is obtained by using $d$-projectives. Furthermore, if $M$ and $N$ are assumed to be finitely generated modules, then  the Auslander-Gruson-Jensen duality would imply that these will be isomorphic to the natural transformations $\nat(\Tor_{k+1}^R(M, -), \Tor_{k+1}^R(N, -))$ of homological functors.
\\ (2)  Assume that $\Q:\bullet\lrt\bullet$ is a quiver. Then the path algebra $R\Q$ is a Gorenstein order in the sense of Auslander \cite{auslander1967functors}. Assume that $\G$ is the category consisting of all lattices over $R\Q$, that is, those $R\Q$-modules $M$ which are finitely generated Gorenstein projective over $R$. As noted in  Example \ref{ex}, $\G$ will be a 1-Frobenius subcategory of $R\Q$-$\Md$.
So for any $X, Y\in\G$ and $i\geq 1$, the following isomorphisms hold
\[\begin{array}{lllll}
\Ext^i_{\G}(X, \syz^iY)/{\p} & \cong \nat(\Ext^{i+1}_{\G}(Y, -), \Ext^{i+1}_{\G}(X, -))\\
& \cong \nat(\Ext^{i+1}_{\G}(-, X), \Ext^{i+1}_{\G}(-, Y)).
\end{array}\]
These will be also isomorphic to the group $\nat(\Tor_{i+1}^{R\Q}(X, -), \Tor_{i+1}^{R\Q}(Y, -))$, whenever $X, Y$ are finitely generated modules.\\ (3) Assume that $R$ is complete and $d>0$. As mentioned in Example \ref{ex}(3), the category $\C$ of all artinian $R$-modules is a $d$-Frobenius category, without having projective objects. Then for any pair of objects $M, N\in\C$ and $k\geq d$, one has the following isomorphisms:  
\[\begin{array}{lllll}
\Ext^k_{\C}(M, \syz^kN)/{\p} & \cong \nat(\Ext^{k+1}_{\C}(N, -), \Ext^{k+1}_{\C}(M, -))\\
& \cong \nat(\Ext^{k+1}_{\C}(-, M), \Ext^{k+1}_{\C}(-, N)).
\end{array}\]

\end{example}

%{\bf{Acknowledgements.}} 

%%%%%%%%%%%%%%%%%%%%%%%%%%%%%%%%%%%%%%%%%%%%%%%%%%%%%%%
\begin{thebibliography}{10}

\bibitem{aus}
{\sc M. Auslander}, Coherent functors,  Proceedings of the Conference of Categorical
Algebra, La Jolla 1965, Springer-Verlag, 1966.

\bibitem{auslander1967functors}
{\sc M.~Auslander}, {\em Functors and morphisms determined by objects}, Representation theory of algebras
(Proc. Conf., Temple Univ., Philadelphia, Pa., 1976),  1–244. Lecture Notes in Pure Appl.
Math., Vol. 37, Dekker, New York, 1978.

\bibitem{aus1}
{\sc M. Auslander}, Isolated singularities and almost split sequences, in: Representation Theory II, in: Lecture Notes in Math., vol. 117, Springer-Verlag, 1986, pp. 194-242.

%\bibitem{as}
%{\sc M. Auslander and O. Solberg,} {Relative homology and representation theory I, relative homology and homologically finite categories,} {\em Comm. Alg.} {\bf 21} (9) (1993), 2995--3031.

\bibitem{bfss}
{\sc A. Bahlekeh, F. S. Fotouhi,  Sh. Salarian and A. SartiPzadeh}, Phantom stable category of $n$-Frobenius categories, to appear in Kyoto J. Math.

\bibitem{bfss1}
{\sc A. Bahlekeh, F. S. Fotouhi,  Sh. Salarian and A. SartiPzadeh}, On the  phantom stable category of $n$-Frobenius categories, available at arXiv:2503.12429.

%\bibitem{bs}
%{\sc A. Bahlekeh and Sh. Salarian,} On the vanishing of $\Ext$ and $\Tor$, available at arXiv:2507.05825.

%\bibitem{be}
%{\sc D. Benson,} Phantom maps and modular representation theory III, {\em J. Algebra} {\bf 248} (2002) 747-754.

%\bibitem{be1}
%{\sc D. Benson and  Ph. G. Gnacadja,} Phantom maps and modular representation theory I, {\em Fund. Math.} {\bf 161} (1999) 37-91.

%\bibitem{be2}
%{\sc D. Benson and Ph. G. Gnacadja,} Phantom maps and modular representation theory II,{\em  Algebr. Represent. Theory} {\bf 4} (2001) 395-404.


%\bibitem{bergh2015gorenstein}
%{\sc P.~A. Bergh, S.~Oppermann, and D.~A. J$\phi$rgensen}, {\em The {G}orensteindefect category}, The Quarterly J.  Math. {\bf 66} (2015), 459--471.

%\bibitem{bu}
%{\sc R.-O. Buchweitz}, Maximal Cohen--Macaulay modules and Tate-cohomology over Gorenstein rings, unpublished paper (1986), \texttt{http://hdl.handle.net/1807/16682}.

\bibitem{buh}
{\sc T. Buhler}, Exact categories, {\em Expo. Math.} {\bf 28}(1) (2010), 1-69.

\bibitem{ej} {\sc E. Enochs and O. M. G. Jenda, }{\sl Relative homological algebra, }
Walter de Gruyter, Berlin-New York, 2000

\bibitem{ee}
{\sc E. Enochs, S. Estrada, J. R. García-Rozas and L. Oyonarte}, Flat and cotorsion quasi-coherent sheaves. Applications,
{\em Algebras and representation theory} {\bf 7} (2004), 441-456.

%\bibitem{gn}
%{\sc Ph.G. Gnacadja,} Phantom maps in the stable module category, {\em J. Algebra} {\bf 201} (2) (1998) 686-702.

\bibitem{gj}
{\sc L. Gruson and C. U. Jensen,}  Dimensions cohomologiques reli$\acute{e}$es aux
foncteurs $\underleftarrow{lim}^{(i)}$, Lecture Notes in Mathematics 867, Springer-Verlag, 1981.

\bibitem{gu}{\sc V. Gulisz}, A functorial approach to $n$-abelian categories, {\em   Comptes Rendus. Mathématique}, {\bf 363} (2025), 1123-1175.

%\bibitem{ha1}
%{\sc D. Happel,} Triangulated categories in the representation theory of finite dimensional algebras, Cambridge University Press,1988.


%\bibitem{he}
%{\sc I. Herzog,} The phantom cover of a module, {\em Adv. Math.} {\bf 215} (1) (2007) 220-249.

%\bibitem{hext}
%{\sc I. Herzog,} Contravariant functors on the category of finitely presented modules, {\em Israel J. Math.} {\bf 167}  (2008), 347-410.

\bibitem{hr}
{\sc P. J. Hilton and D. Rees,} Natural maps of extension functors and a theorem of R. G. Swan, {\em Proc. Cambridge
Philos. Soc.} {\bf 57} (1961), 489-502.

\bibitem{hs}
{\sc P. J. Hilton and U. Stammbach,} A course in homological algebra. chapter IV. Derived Functors,
pages 116-165. Springer-Verlag, 1971.

%\bibitem{hz}
%{\sc C. Holdaway and K. Zatloukal}, The stable category of a Frobenius category is triangulated, Unpublished lecture notes (2012). 

\bibitem{iy}
{\sc O. Iyama,} Higher-dimensional Auslander-Reiten theory on maximal orthogonal subcategories,
{\em Adv. Math.} {\bf 210}  (2007), 22-50.

%\bibitem{ke}
%{\sc B. Keller,} On differential graded categories, in: {\em International Congress of Mathematicians}. Vol. II, European Mathematical Society, Zürich (2006), 151–190.

%\bibitem{lz}
%{\sc Y. Liu and P. Zhou,}  Frobenius $n$-exangulated categories, {\em J. Algebra} {\bf 559} (2020), 161-183.

%\bibitem{mao}
%{\sc L. Mao,} Higher phantom and $\Ext$-phantom morphisms, {\em J. Algebra Appl.} {\bf 17} (2018), no. 1, 1850012, 15 pp.

%\bibitem{mao1}{\sc L. Mao,} Higher phantom  morphisms with respect to a subfunctor of $\Ext$, {\em Algebr Represent Theor} {\bf 22} (2019), 407-424.

\bibitem{mart}{\sc A. Martsinkovsky}, On Direct Summands of Homological Functors on Length Categories, {\em  Appl. Cat. Str.}, {\bf 24}(4) (2016), 421-431.

%\bibitem{mc}
%{\sc C. A. McGibbon,} Phantom maps, in: Handbook of Algebraic Topology, Elsevier Science B.V., Amsterdam, 1995, pp. 1209-1257.

\bibitem{mit}
{\sc B. Mitchel,}  Theory of categories, Pure and applied Math. A series of monographs and textbooks, vol. 17, Elsevier, Academic Press, New York, London, 1965.

%\bibitem{mu}
%{\sc D. Murfet,}  The  mock  homotopy  category  of  projectives  and Grothendieck  duality,  PhD  thesis, available online at http://www.therisingsea.org/thesis.pdf, 2007.

%\bibitem{ms}
%{\sc D. Murfet and Sh. Salarian,}  Totally acyclic complexes over noetherian schemes, {\em Adv. Math.} {\bf 226} (2011), 1096-1133.


%\bibitem{ne}
%{\sc A. Neeman,} The Brown representability theorem and phantomless triangulated categories, {\em J. Algebra} {\bf 151} (1) (1992), 118-155.

%\bibitem{ne2}
%{\sc A. Neeman,}  The homotopy category of flat modules, and Grothendieck duality, {\em Invent. Math.} {\bf 174} (2008),  255-308.

\bibitem{rot}
{\sc J. J. Rotman, } An introduction to homological algebra, 
New York, Springer, 2009.

%\bibitem{ver}
%{\sc J. L. Verdier,}  Des cat$\acute{e}$gories d$\acute{e}$riv$\acute{e}$es des cat$\acute{e}$gories ab$\acute{e}$liennes,  {\em Ast$\acute{e}$risque} No. 239 (1996).


\bibitem{yo}{\sc Y. Yoshino,}  Cohen-Macaulay Modules over Cohen-Macaulay Rings,
Lond. Math. Soc. Lecture Notes Series, 146, Cambridge Univ. Press, Cambridge, 1990.

%\bibitem{xu}
%{\sc Y. Xu}, Flat covers of modules, Lecture Notes in Mathematics, vol. 1634 (Springer, 1996).

\end{thebibliography}
%%%%%%%%%%%%%%%%%%%%%%%%%%%%%%%%%%%%%%%%%%%%%%%%%%%%%%%%

\end{document}

